% Options for packages loaded elsewhere
\PassOptionsToPackage{unicode}{hyperref}
\PassOptionsToPackage{hyphens}{url}
\PassOptionsToPackage{dvipsnames,svgnames,x11names}{xcolor}
%
\documentclass[
  10pt,
]{article}
\usepackage{amsmath,amssymb}
\usepackage{iftex}
\ifPDFTeX
  \usepackage[T1]{fontenc}
  \usepackage[utf8]{inputenc}
  \usepackage{textcomp} % provide euro and other symbols
\else % if luatex or xetex
  \usepackage{unicode-math} % this also loads fontspec
  \defaultfontfeatures{Scale=MatchLowercase}
  \defaultfontfeatures[\rmfamily]{Ligatures=TeX,Scale=1}
\fi
\usepackage{lmodern}
\ifPDFTeX\else
  % xetex/luatex font selection
\fi
% Use upquote if available, for straight quotes in verbatim environments
\IfFileExists{upquote.sty}{\usepackage{upquote}}{}
\IfFileExists{microtype.sty}{% use microtype if available
  \usepackage[]{microtype}
  \UseMicrotypeSet[protrusion]{basicmath} % disable protrusion for tt fonts
}{}
\makeatletter
\@ifundefined{KOMAClassName}{% if non-KOMA class
  \IfFileExists{parskip.sty}{%
    \usepackage{parskip}
  }{% else
    \setlength{\parindent}{0pt}
    \setlength{\parskip}{6pt plus 2pt minus 1pt}}
}{% if KOMA class
  \KOMAoptions{parskip=half}}
\makeatother
\usepackage{xcolor}
\usepackage[left=2cm, right=2cm, top=2cm, bottom=3cm, footskip = .5cm]{geometry}
\usepackage{longtable,booktabs,array}
\usepackage{calc} % for calculating minipage widths
% Correct order of tables after \paragraph or \subparagraph
\usepackage{etoolbox}
\makeatletter
\patchcmd\longtable{\par}{\if@noskipsec\mbox{}\fi\par}{}{}
\makeatother
% Allow footnotes in longtable head/foot
\IfFileExists{footnotehyper.sty}{\usepackage{footnotehyper}}{\usepackage{footnote}}
\makesavenoteenv{longtable}
\usepackage{graphicx}
\makeatletter
\newsavebox\pandoc@box
\newcommand*\pandocbounded[1]{% scales image to fit in text height/width
  \sbox\pandoc@box{#1}%
  \Gscale@div\@tempa{\textheight}{\dimexpr\ht\pandoc@box+\dp\pandoc@box\relax}%
  \Gscale@div\@tempb{\linewidth}{\wd\pandoc@box}%
  \ifdim\@tempb\p@<\@tempa\p@\let\@tempa\@tempb\fi% select the smaller of both
  \ifdim\@tempa\p@<\p@\scalebox{\@tempa}{\usebox\pandoc@box}%
  \else\usebox{\pandoc@box}%
  \fi%
}
% Set default figure placement to htbp
\def\fps@figure{htbp}
\makeatother
\setlength{\emergencystretch}{3em} % prevent overfull lines
\providecommand{\tightlist}{%
  \setlength{\itemsep}{0pt}\setlength{\parskip}{0pt}}
\setcounter{secnumdepth}{-\maxdimen} % remove section numbering
% Set up the fonts
\usepackage[urw-palatino]{mathdesign}
\usepackage[T1]{fontenc}

% Add accessibility support from http://www.richschwinn.com/accessibility
\RequirePackage{accsupp}
\RequirePackage{pdfcomment}
\newcommand{\AccTool}[2]{\BeginAccSupp{method=pdfstringdef,unicode,Alt={{#1}}}\pdftooltip{{#2}}{{#1}}\EndAccSupp{}}

% Set the language for 508
\hypersetup{
  pdftitle = {title},
  pdflang = en-US}

% left-align all figure captions 
\usepackage{caption}
\captionsetup{justification=raggedright,singlelinecheck=false}

% Set up the headers and footers
\usepackage{graphicx}
\usepackage{fancyhdr}
\usepackage{ifthen}
%\usepackage{everypage-1x}
\usepackage{float}
%\usepackage{subfig}
%\usepackage{subcaption}

% Avoid struggling over figure and table float in Rmarkdown
\let\origfigure\figure
\let\endorigfigure\endfigure
\renewenvironment{figure}[1][2] {
    \expandafter\origfigure\expandafter[H]
} {
    \endorigfigure
}

\let\origtable\table
\let\endorigtable\endtable
\renewenvironment{table}[1][2] {
    \expandafter\origtable\expandafter[H]
} {
    \endorigtable
}

% First page has the large title and NOAA logo
\pagestyle{fancy}
\fancyhf{}
\setlength\headheight{40pt}
\fancyheadoffset[L]{0.5cm}
\cfoot{\thepage}

\fancyheadinit{%
   \ifthenelse{\value{page}=5}%
      {\fancyhead[R]{
\includegraphics[width=40pt]{images/NOAA_logo.png} \\ \textsf{\emph{Feb 13, 2026}}}
       \fancyhead[L]{\textsf{\LARGE State of the Ecosystem 2025: REVIEW DRAFT}}
      }%
      {\fancyhead[R]{}
       \fancyhead[L]{\textsf{\emph{State of the Ecosystem 2025: REVIEW DRAFT}}}
      }
}



\renewcommand{\headrulewidth}{0.4pt}
\renewcommand{\footrulewidth}{0pt}

% Make caption fonts a bit smaller
\usepackage[font={small}]{caption}


% Change section labels to san serif
\usepackage{sectsty}
\allsectionsfont{\normalfont\sffamily\bfseries}
\usepackage{multirow}
\usepackage{multicol}
\usepackage{colortbl}
\usepackage{hhline}
\newlength\Oldarrayrulewidth
\newlength\Oldtabcolsep
\usepackage{longtable}
\usepackage{array}
\usepackage{hyperref}
\usepackage{float}
\usepackage{wrapfig}
\usepackage{bookmark}
\IfFileExists{xurl.sty}{\usepackage{xurl}}{} % add URL line breaks if available
\urlstyle{same}
\hypersetup{
  colorlinks=true,
  linkcolor={Maroon},
  filecolor={Maroon},
  citecolor={Blue},
  urlcolor={blue},
  pdfcreator={LaTeX via pandoc}}

\author{}
\date{\vspace{-2.5em}}

\begin{document}

\setcounter{page}{4}
\thispagestyle{fancy}

\section{Introduction}\label{introduction}

\subsection{About This Report}\label{about-this-report}

This report is for the New England Fishery Management Council (NEFMC). The purpose of this report is to synthesize ecosystem information to allow the NEFMC to better meet fishery management objectives. The major messages of the report are synthesized on pages 1-3, with highlights of 2025 ecosystem events on page 4.

The information in this report is organized into two main sections; \hyperref[performance-relative-to-fishery-management-objectives]{performance measured against ecosystem-level management objectives} (Table \ref{tab:management-objectives}), and potential \hyperref[risks-to-meeting-fishery-management-objectives]{risks to meeting fishery management objectives} (Table \ref{tab:management-risks}: \hyperref[climate-and-ecosystem-change]{climate change} and \hyperref[other-ocean-uses-offshore-wind]{other ocean uses}). A final section highlights \hyperref[highlights]{notable 2025 ecosystem observations}.

\subsection{Report structure}\label{report-structure}

A glossary of terms\footnote{\url{https://noaa-edab.github.io/tech-doc/glossary.html}}, detailed technical methods documentation\footnote{\url{https://noaa-edab.github.io/tech-doc/}}, indicator data\footnote{\url{https://noaa-edab.github.io/ecodata/}}, and detailed indicator descriptions\footnote{\url{https://noaa-edab.github.io/catalog/index.html}} are available online. We recommend new readers first review the details of standard figure formatting (Fig. \ref{fig:docformat}a), categorization of fish and invertebrate species into feeding guilds (Table \ref{tab:species-groupings}), and definitions of ecological production units (EPUs, including the Gulf of Maine (GOM) and Georges Bank (GB); Fig. \ref{fig:docformat}b) provided at the end of the document.

The two main sections contain subsections for each management objective or potential risk. Within each subsection, we first review observed trends for indicators representing each objective or risk, including the status of the most recent data year relative to a threshold (if available) or relative to the long-term average. Second, we identify potential drivers of observed trends, and synthesize results of indicators related to those drivers to outline potential implications for management. For example, if there are multiple drivers related to an indicator trend, do indicators associated with the drivers have similar trends, and can any drivers be affected by management action(s)? We emphasize that these implications are intended to represent testable hypotheses at present, rather than ``answers,'' because the science behind these indicators and syntheses continues to develop.

\global\setlength{\Oldarrayrulewidth}{\arrayrulewidth}

\global\setlength{\Oldtabcolsep}{\tabcolsep}

\setlength{\tabcolsep}{2pt}

\renewcommand*{\arraystretch}{1.5}



\providecommand{\ascline}[3]{\noalign{\global\arrayrulewidth #1}\arrayrulecolor[HTML]{#2}\cline{#3}}

\begin{longtable}[c]{|p{1.77in}|p{4.47in}}

\caption{Ecosystem-scale\ fishery\ management\ objectives\ in\ New\ England}\label{tab:management-objectives}\\

\ascline{1.5pt}{666666}{1-2}

\multicolumn{1}{>{\raggedright}m{\dimexpr 1.77in+0\tabcolsep}}{\textcolor[HTML]{000000}{\fontsize{9}{9}\selectfont{Objective\ categories}}} & \multicolumn{1}{>{\raggedright}m{\dimexpr 4.47in+0\tabcolsep}}{\textcolor[HTML]{000000}{\fontsize{9}{9}\selectfont{Indicators\ reported}}} \\

\ascline{1.5pt}{666666}{1-2}\endfirsthead \caption[]{Ecosystem-scale\ fishery\ management\ objectives\ in\ New\ England}\label{tab:management-objectives}\\

\ascline{1.5pt}{666666}{1-2}

\multicolumn{1}{>{\raggedright}m{\dimexpr 1.77in+0\tabcolsep}}{\textcolor[HTML]{000000}{\fontsize{9}{9}\selectfont{Objective\ categories}}} & \multicolumn{1}{>{\raggedright}m{\dimexpr 4.47in+0\tabcolsep}}{\textcolor[HTML]{000000}{\fontsize{9}{9}\selectfont{Indicators\ reported}}} \\

\ascline{1.5pt}{666666}{1-2}\endhead



\multicolumn{2}{>{\raggedright}m{\dimexpr 6.24in+2\tabcolsep}}{\textcolor[HTML]{000000}{\fontsize{9}{9}\selectfont{\textbf{Objectives:\ Provisioning\ and\ Cultural\ Services}}}} \\





\multicolumn{1}{>{\raggedright}m{\dimexpr 1.77in+0\tabcolsep}}{\textcolor[HTML]{000000}{\fontsize{9}{9}\selectfont{Seafood\ Production}}} & \multicolumn{1}{>{\raggedright}m{\dimexpr 4.47in+0\tabcolsep}}{\textcolor[HTML]{000000}{\fontsize{9}{9}\selectfont{Landings;\ commercial\ total\ and\ by\ feeding\ guild;\ recreational\ harvest}}} \\





\multicolumn{1}{>{\raggedright}m{\dimexpr 1.77in+0\tabcolsep}}{\textcolor[HTML]{000000}{\fontsize{9}{9}\selectfont{Commercial\ Profits}}} & \multicolumn{1}{>{\raggedright}m{\dimexpr 4.47in+0\tabcolsep}}{\textcolor[HTML]{000000}{\fontsize{9}{9}\selectfont{Revenue\ decomposed\ to\ price\ and\ volume}}} \\





\multicolumn{1}{>{\raggedright}m{\dimexpr 1.77in+0\tabcolsep}}{\textcolor[HTML]{000000}{\fontsize{9}{9}\selectfont{Recreational\ Opportunities}}} & \multicolumn{1}{>{\raggedright}m{\dimexpr 4.47in+0\tabcolsep}}{\textcolor[HTML]{000000}{\fontsize{9}{9}\selectfont{Angler\ trips;\ recreational\ fleet\ diversity}}} \\





\multicolumn{1}{>{\raggedright}m{\dimexpr 1.77in+0\tabcolsep}}{\textcolor[HTML]{000000}{\fontsize{9}{9}\selectfont{Stability}}} & \multicolumn{1}{>{\raggedright}m{\dimexpr 4.47in+0\tabcolsep}}{\textcolor[HTML]{000000}{\fontsize{9}{9}\selectfont{Fishery\ and\ ecosystem\ volatility,\ adaptive\ capacity,\ and\ shifts\ from\ baseline}}} \\





\multicolumn{1}{>{\raggedright}m{\dimexpr 1.77in+0\tabcolsep}}{\textcolor[HTML]{000000}{\fontsize{9}{9}\selectfont{Social\ \&\ Cultural}}} & \multicolumn{1}{>{\raggedright}m{\dimexpr 4.47in+0\tabcolsep}}{\textcolor[HTML]{000000}{\fontsize{9}{9}\selectfont{Community\ fishing\ engagement\ and\ social\ vulnerability\ status}}} \\





\multicolumn{1}{>{\raggedright}m{\dimexpr 1.77in+0\tabcolsep}}{\textcolor[HTML]{000000}{\fontsize{9}{9}\selectfont{Protected\ Species}}} & \multicolumn{1}{>{\raggedright}m{\dimexpr 4.47in+0\tabcolsep}}{\textcolor[HTML]{000000}{\fontsize{9}{9}\selectfont{Bycatch;\ population\ (adult\ and\ juvenile)\ numbers;\ mortalities}}} \\





\multicolumn{2}{>{\raggedright}m{\dimexpr 6.24in+2\tabcolsep}}{\textcolor[HTML]{000000}{\fontsize{9}{9}\selectfont{\textbf{Potential\ Drivers:\ Supporting\ and\ Regulating\ Services}}}} \\





\multicolumn{1}{>{\raggedright}m{\dimexpr 1.77in+0\tabcolsep}}{\textcolor[HTML]{000000}{\fontsize{9}{9}\selectfont{Management}}} & \multicolumn{1}{>{\raggedright}m{\dimexpr 4.47in+0\tabcolsep}}{\textcolor[HTML]{000000}{\fontsize{9}{9}\selectfont{Stock\ status;\ catch\ compared\ with\ catch\ limits}}} \\





\multicolumn{1}{>{\raggedright}m{\dimexpr 1.77in+0\tabcolsep}}{\textcolor[HTML]{000000}{\fontsize{9}{9}\selectfont{Biomass}}} & \multicolumn{1}{>{\raggedright}m{\dimexpr 4.47in+0\tabcolsep}}{\textcolor[HTML]{000000}{\fontsize{9}{9}\selectfont{Biomass\ or\ abundance\ by\ feeding\ guild\ from\ surveys}}} \\





\multicolumn{1}{>{\raggedright}m{\dimexpr 1.77in+0\tabcolsep}}{\textcolor[HTML]{000000}{\fontsize{9}{9}\selectfont{Environment}}} & \multicolumn{1}{>{\raggedright}m{\dimexpr 4.47in+0\tabcolsep}}{\textcolor[HTML]{000000}{\fontsize{9}{9}\selectfont{Climate\ and\ ecosystem\ risk\ indicators\ listed\ in\ Table\ 2}}} \\

\ascline{1.5pt}{666666}{1-2}



\end{longtable}



\arrayrulecolor[HTML]{000000}

\global\setlength{\arrayrulewidth}{\Oldarrayrulewidth}

\global\setlength{\tabcolsep}{\Oldtabcolsep}

\renewcommand*{\arraystretch}{1}

\newpage

\global\setlength{\Oldarrayrulewidth}{\arrayrulewidth}

\global\setlength{\Oldtabcolsep}{\tabcolsep}

\setlength{\tabcolsep}{2pt}

\renewcommand*{\arraystretch}{1.5}



\providecommand{\ascline}[3]{\noalign{\global\arrayrulewidth #1}\arrayrulecolor[HTML]{#2}\cline{#3}}

\begin{longtable}[c]{|p{1.00in}|p{2.20in}|p{2.80in}}

\caption{Risks\ to\ meeting\ fishery\ management\ objectives\ in\ New\ England}\label{tab:management-risks}\\

\ascline{1.5pt}{666666}{1-3}

\multicolumn{1}{>{\raggedright}m{\dimexpr 1in+0\tabcolsep}}{\textcolor[HTML]{000000}{\fontsize{9}{9}\selectfont{Risk\ categories}}} & \multicolumn{1}{>{\raggedright}m{\dimexpr 2.2in+0\tabcolsep}}{\textcolor[HTML]{000000}{\fontsize{9}{9}\selectfont{Observation\ indicators\ reported}}} & \multicolumn{1}{>{\raggedright}m{\dimexpr 2.8in+0\tabcolsep}}{\textcolor[HTML]{000000}{\fontsize{9}{9}\selectfont{Potential\ driver\ indicators\ reported}}} \\

\ascline{1.5pt}{666666}{1-3}\endfirsthead \caption[]{Risks\ to\ meeting\ fishery\ management\ objectives\ in\ New\ England}\label{tab:management-risks}\\

\ascline{1.5pt}{666666}{1-3}

\multicolumn{1}{>{\raggedright}m{\dimexpr 1in+0\tabcolsep}}{\textcolor[HTML]{000000}{\fontsize{9}{9}\selectfont{Risk\ categories}}} & \multicolumn{1}{>{\raggedright}m{\dimexpr 2.2in+0\tabcolsep}}{\textcolor[HTML]{000000}{\fontsize{9}{9}\selectfont{Observation\ indicators\ reported}}} & \multicolumn{1}{>{\raggedright}m{\dimexpr 2.8in+0\tabcolsep}}{\textcolor[HTML]{000000}{\fontsize{9}{9}\selectfont{Potential\ driver\ indicators\ reported}}} \\

\ascline{1.5pt}{666666}{1-3}\endhead



\multicolumn{3}{>{\raggedright}m{\dimexpr 6in+4\tabcolsep}}{\textcolor[HTML]{000000}{\fontsize{9}{9}\selectfont{\textbf{Climate\ and\ Ecosystem\ Risks}}}} \\





\multicolumn{1}{>{\raggedright}m{\dimexpr 1in+0\tabcolsep}}{\textcolor[HTML]{000000}{\fontsize{9}{9}\selectfont{Risks\ to\ Managing\ Spatially}}} & \multicolumn{1}{>{\raggedright}m{\dimexpr 2.2in+0\tabcolsep}}{\textcolor[HTML]{000000}{\fontsize{9}{9}\selectfont{Managed\ species\ (fish\ and\ cetacean)\ distribution\ shifts}}} & \multicolumn{1}{>{\raggedright}m{\dimexpr 2.8in+0\tabcolsep}}{\textcolor[HTML]{000000}{\fontsize{9}{9}\selectfont{Benthic\ and\ pelagic\ forage\ distribution;\ ocean\ temperature,\ changes\ in\ currents\ and\ cold\ pool}}} \\





\multicolumn{1}{>{\raggedright}m{\dimexpr 1in+0\tabcolsep}}{\textcolor[HTML]{000000}{\fontsize{9}{9}\selectfont{Risks\ to\ Managing\ Seasonally}}} & \multicolumn{1}{>{\raggedright}m{\dimexpr 2.2in+0\tabcolsep}}{\textcolor[HTML]{000000}{\fontsize{9}{9}\selectfont{Managed\ species\ spawning\ and\ migration\ timing\ changes}}} & \multicolumn{1}{>{\raggedright}m{\dimexpr 2.8in+0\tabcolsep}}{\textcolor[HTML]{000000}{\fontsize{9}{9}\selectfont{Habitat\ timing:\ Length\ of\ ocean\ summer,\ cold\ pool\ seasonal\ persistence}}} \\





\multicolumn{1}{>{\raggedright}m{\dimexpr 1in+0\tabcolsep}}{\textcolor[HTML]{000000}{\fontsize{9}{9}\selectfont{Risks\ to\ Setting\ Catch\ Limits}}} & \multicolumn{1}{>{\raggedright}m{\dimexpr 2.2in+0\tabcolsep}}{\textcolor[HTML]{000000}{\fontsize{9}{9}\selectfont{Managed\ species\ body\ condition\ and\ recruitment\ changes}}} & \multicolumn{1}{>{\raggedright}m{\dimexpr 2.8in+0\tabcolsep}}{\textcolor[HTML]{000000}{\fontsize{9}{9}\selectfont{Benthic\ and\ pelagic\ forage\ quality\ \&\ abundance:\ ocean\ temperature\ \&\ acidification\ }}} \\





\multicolumn{3}{>{\raggedright}m{\dimexpr 6in+4\tabcolsep}}{\textcolor[HTML]{000000}{\fontsize{9}{9}\selectfont{\textbf{Other\ Ocean\ Uses\ Risks}}}} \\





\multicolumn{1}{>{\raggedright}m{\dimexpr 1in+0\tabcolsep}}{\textcolor[HTML]{000000}{\fontsize{9}{9}\selectfont{Offshore\ Wind\ Risks}}} & \multicolumn{1}{>{\raggedright}m{\dimexpr 2.2in+0\tabcolsep}}{\textcolor[HTML]{000000}{\fontsize{9}{9}\selectfont{Fishery\ revenue\ and\ landings\ from\ wind\ lease\ areas\ by\ species\ and\ port}}} & \multicolumn{1}{>{\raggedright}m{\dimexpr 2.8in+0\tabcolsep}}{\textcolor[HTML]{000000}{\fontsize{9}{9}\selectfont{Wind\ development\ speed;\ Protected\ species\ presence\ and\ \ hotspots}}} \\

\ascline{1.5pt}{666666}{1-3}



\end{longtable}



\arrayrulecolor[HTML]{000000}

\global\setlength{\arrayrulewidth}{\Oldarrayrulewidth}

\global\setlength{\tabcolsep}{\Oldtabcolsep}

\renewcommand*{\arraystretch}{1}

\newpage

\section{Performance Relative to Fishery Management Objectives}\label{performance-relative-to-fishery-management-objectives}

In this section, we examine indicators related to broad, ecosystem-level fishery management objectives. We also provide hypotheses on the implications of these trends---why we are seeing them, what's driving them, and potential or observed regime shifts or changes in ecosystem structure. Identifying multiple drivers, regime shifts, and potential changes to ecosystem structure, as well as identifying the most vulnerable resources, can help managers determine whether anything needs to be done differently to meet objectives and how to prioritize upcoming issues/risks.

\subsection{Seafood Production}\label{seafood-production}

\subsubsection{Indicators: Landings; commercial and recreational}\label{indicators-landings-commercial-and-recreational}

In New England, total \href{https://noaa-edab.github.io/catalog/comdat.html}{commercial landings} (includes bait and industrial uses), U.S. seafood landings (excludes industrial and bait uses), and NEFMC-managed landings have long-term declines (Fig. 1). Declines are seen in commercial landings by \href{https://noaa-edab.github.io/catalog/species_groupings.html}{guild}, which include all species and all uses caught with GB and the GOM, and are reported in total and for NEFMC managed species only. Downward trends persist for piscivores and benthivores in both regions. Current high total landings for benthivores (GOM) are attributable to American lobster. High benthos landings (GB) are attributable to clams and scallops, although they are below the long-term mean in 2024 (Fig. \ref{fig:total-landings}). Current landings of planktivores and piscivores remain among the lowest points in the time series.

\begin{figure}

{\centering \includegraphics[width=1\linewidth]{../../../AppData/Local/Temp/RtmpIjNXEs/file2644409e58ba} 

}

\caption{Total commercial landings (black), total U.S. seafood landings (blue), and New England managed U.S. seafood landings (red) for Georges Bank (GB, top) and the Gulf of Maine (GOM, bottom),  with significant decline (purple) in total landings.}\label{fig:total-landings}
\end{figure}

\begin{figure}

{\centering \includegraphics[width=1\linewidth]{../../../AppData/Local/Temp/RtmpIjNXEs/file264426f479fd} 

}

\caption{Total commercial landings (black) and NEFMC managed U.S seafood landings (red) by feeding guild for the Gulf of Maine (GOM, right) and Georges Bank (GB, left), with significant long-term declines (purple) in all categories.}\label{fig:comm-landings}
\end{figure}

New England ports face a moderate risk from environmental variability with no long-term trend, as evaluated by the 2025 \href{https://noaa-edab.github.io/catalog/community_climate_vulnerability.html}{Community Environmental Variability Risk Indicators}indicators (Fig. \ref{fig:climatevul-land}). These indicators assess port level risk to environmental variability based on dependence on species and their respective bioenvironmental vulnerabilities as assessed by regional experts. Total Vulnerability measures how much of a region's landings (or revenue) is dependent on species that are sensitive to different climate and environmental change factors including temperature and acidification.

\begin{figure}

{\centering \includegraphics[width=1\linewidth]{../../../AppData/Local/Temp/RtmpIjNXEs/file26441855ca5} 

}

\caption{New England region total climate vulnerability of commercial landings (sum of New England port landings weighted by species climate vulnerability from Hare et al. 2016). Horizontal colored bars show different climate risk levels.}\label{fig:climatevul-land}
\end{figure}

In New England, \href{https://noaa-edab.github.io/catalog/recdat.html}{recreational harvest} (retained fish presumed to be eaten) shows a long-term decline with recent harvest remaining at a time series low (Fig. \ref{fig:rec-landings}). This pattern may indicate a shift towards catch-and-release strategies as opposed to catch for harvest. \href{https://noaa-edab.github.io/catalog/rec_hms.html}{Recreational shark landings} have generally decreased for most shark groups through 2024 (Fig. \ref{fig:rec-hms-ne}). The recent low in pelagic shark landings is largely driven by regulatory changes implemented in 2018, followed by the closure of the shortfin mako fishery in 2022. These actions were intended to rebuild the North Atlantic shortfin mako stock and comply with binding recommendations by the International Commission for the Conservation of Atlantic Tunas (ICCAT).

\begin{figure}

{\centering \includegraphics[width=1\linewidth]{../../../AppData/Local/Temp/RtmpIjNXEs/file26442d7941c2} 

}

\caption{Total recreational seafood harvest (millions of pounds, black, significant decrease, purple) in the New England region.}\label{fig:rec-landings}
\end{figure}

\begin{figure}

{\centering \includegraphics[width=1\linewidth]{../../../AppData/Local/Temp/RtmpIjNXEs/file264470217fa3} 

}

\caption{Recreational shark landings in the New England region from NOAA Fisheries' Large Pelagics Survey.}\label{fig:rec-hms-ne}
\end{figure}

\href{https://noaa-edab.github.io/catalog/aquaculture.html?q=aqua\#aquaculture}{Aquaculture} can comprise a significant proportion of seafood production in specific communities, but not all aquaculture production is included in total seafood landings above. In 2022, the Northeast region produced approximately 6,300 metric tons of aquacultured shellfish, with revenue of \$133 million (Fisheries of the United States, 2022).

\subsubsection{Implications}\label{implications}

Declining commercial landings (total and seafood) and recreational harvest can be attributed to many interacting factors, including combinations of ecosystem and stock production, management actions, market conditions, and environmental change. While we cannot evaluate all possible drivers at present, here we evaluate the extent to which stock status, management, and system biomass trends may play a role.

\paragraph{Stock Status}\label{stock-status}

Single species \href{https://noaa-edab.github.io/catalog/stock_status.html}{management objectives} (1. maintaining biomass above minimum thresholds and 2. maintaining fishing mortality below overfishing limits) are not being met for some NEFMC managed species. Specifically, 17 stocks are currently estimated to be below B\textsubscript{MSY} targets and 10 below B\textsubscript{MSY} thresholds (Fig. \ref{fig:stock-status}) However, the status of 12 stocks is unknown (Table \ref{tab:stock-status-table}). Although stock status and associated management constraints are likely contributing to decreased landings, these management actions are enacted in response to biomass, where less conservative regulations would not necessarily mean higher landings. To better address the role of management in future reports, we could examine how the total allowable catch (TAC) and the percentage of the TAC utilized for each species has changed through time.

\begin{figure}

{\centering \includegraphics[width=1\linewidth]{../../../AppData/Local/Temp/RtmpIjNXEs/file26442c05127f} 

}

\caption{Summary of single species status for NEFMC and jointly federally managed stocks (goosefish and spiny dogfish).  The dotted vertical line at one is the target biomass reference point of B.  The dashed lines are the management thresholds of B (vertical) or F (horizontal). Stocks with a $B/B_{MSY}$ estimate but without an $F/F_{MSY}$ estimate are denoted in a separate box plot (top). Colors denote stocks with $B/B_{MSY}$ < 0.5 or $F/F_{MSY}$ (orange), stocks 0.5<$B/B_{MSY}$<1 (blue), and stocks $B/B_{MSY}$>1 (green).CCGOM = Cape Cod Gulf of Maine, GOM = Gulf of Maine, GB = Georges Bank, SNEMA = Southern New England Mid Atlantic}\label{fig:stock-status}
\end{figure}

\newpage

\global\setlength{\Oldarrayrulewidth}{\arrayrulewidth}

\global\setlength{\Oldtabcolsep}{\tabcolsep}

\setlength{\tabcolsep}{2pt}

\renewcommand*{\arraystretch}{1}



\providecommand{\ascline}[3]{\noalign{\global\arrayrulewidth #1}\arrayrulecolor[HTML]{#2}\cline{#3}}

\begin{longtable}[c]{|p{3.13in}}

\caption{Unknown\ or\ partially\ known\ stock\ status\ for\ NEFMC\ and\ jointly\ managed\ species.}\label{tab:stock-status-table}\\

\ascline{1.5pt}{666666}{1-1}

\multicolumn{1}{>{\raggedright}m{\dimexpr 3.13in+0\tabcolsep}}{\textcolor[HTML]{000000}{\fontsize{9}{9}\selectfont{Stock}}} \\

\ascline{1.5pt}{666666}{1-1}\endfirsthead \caption[]{Unknown\ or\ partially\ known\ stock\ status\ for\ NEFMC\ and\ jointly\ managed\ species.}\label{tab:stock-status-table}\\

\ascline{1.5pt}{666666}{1-1}

\multicolumn{1}{>{\raggedright}m{\dimexpr 3.13in+0\tabcolsep}}{\textcolor[HTML]{000000}{\fontsize{9}{9}\selectfont{Stock}}} \\

\ascline{1.5pt}{666666}{1-1}\endhead



\multicolumn{1}{>{\raggedright}m{\dimexpr 3.13in+0\tabcolsep}}{\textcolor[HTML]{000000}{\fontsize{9}{9}\selectfont{Atlantic\ salmon\ -\ Gulf\ of\ Maine}}} \\





\multicolumn{1}{>{\raggedright}m{\dimexpr 3.13in+0\tabcolsep}}{\textcolor[HTML]{000000}{\fontsize{9}{9}\selectfont{Red\ deepsea\ crab\ -\ Northwestern\ Atlantic}}} \\





\multicolumn{1}{>{\raggedright}m{\dimexpr 3.13in+0\tabcolsep}}{\textcolor[HTML]{000000}{\fontsize{9}{9}\selectfont{Atlantic\ cod\ -\ Gulf\ of\ Maine}}} \\





\multicolumn{1}{>{\raggedright}m{\dimexpr 3.13in+0\tabcolsep}}{\textcolor[HTML]{000000}{\fontsize{9}{9}\selectfont{Atlantic\ halibut\ -\ Northwestern\ Atlantic\ Coast}}} \\





\multicolumn{1}{>{\raggedright}m{\dimexpr 3.13in+0\tabcolsep}}{\textcolor[HTML]{000000}{\fontsize{9}{9}\selectfont{Offshore\ hake\ -\ Northwestern\ Atlantic\ Coast}}} \\





\multicolumn{1}{>{\raggedright}m{\dimexpr 3.13in+0\tabcolsep}}{\textcolor[HTML]{000000}{\fontsize{9}{9}\selectfont{Red\ hake\ -\ Gulf\ of\ Maine\ /\ Northern\ Georges\ Bank}}} \\





\multicolumn{1}{>{\raggedright}m{\dimexpr 3.13in+0\tabcolsep}}{\textcolor[HTML]{000000}{\fontsize{9}{9}\selectfont{Red\ hake\ -\ Southern\ Georges\ Bank\ /\ Mid-Atlantic}}} \\





\multicolumn{1}{>{\raggedright}m{\dimexpr 3.13in+0\tabcolsep}}{\textcolor[HTML]{000000}{\fontsize{9}{9}\selectfont{Windowpane\ -\ Gulf\ of\ Maine\ /\ Georges\ Bank}}} \\





\multicolumn{1}{>{\raggedright}m{\dimexpr 3.13in+0\tabcolsep}}{\textcolor[HTML]{000000}{\fontsize{9}{9}\selectfont{Witch\ flounder\ -\ Northwestern\ Atlantic\ Coast}}} \\





\multicolumn{1}{>{\raggedright}m{\dimexpr 3.13in+0\tabcolsep}}{\textcolor[HTML]{000000}{\fontsize{9}{9}\selectfont{Goosefish\ -\ Gulf\ of\ Maine\ /\ Northern\ Georges\ Bank}}} \\





\multicolumn{1}{>{\raggedright}m{\dimexpr 3.13in+0\tabcolsep}}{\textcolor[HTML]{000000}{\fontsize{9}{9}\selectfont{Goosefish\ -\ Southern\ Georges\ Bank\ /\ Mid-Atlantic}}} \\

\ascline{1.5pt}{666666}{1-1}



\end{longtable}



\arrayrulecolor[HTML]{000000}

\global\setlength{\arrayrulewidth}{\Oldarrayrulewidth}

\global\setlength{\tabcolsep}{\Oldtabcolsep}

\renewcommand*{\arraystretch}{1}

\paragraph{System Biomass}\label{system-biomass}

Declining landings are likely driven by the relative abundance of specific targeted species rather than major shifts in ecosystem trophic structure or feeding guilds. Scientific surveys show that \href{https://noaa-edab.github.io/catalog/aggregate_biomass.html}{Aggregate biomass} has been mostly stable or increasing in both regions (Fig. \ref{fig:nefsc-biomass-gb} \& Fig. \ref{fig:nefsc-biomass-gom}). The benthivores biomass recently peaked due to a large haddock recruitment, but appears to be returning to average levels. Planktivore biomass on GB continues to rise due to increased Atlantic mackerel. On GB, trends in piscivores on GB are mixed, and benthos are increasing in both seasons. State-level data show the Massachusetts survey (Fig. \ref{fig:mass-biomass}) mirroring the increase in fall planktivores but noting a spring decrease in fish-eaters and a year-round decline in benthos; the \href{https://noaa-edab.github.io/catalog/mass_inshore_survey.html?q=survey\#mass_inshore_survey}{New Hampshire/Maine survey} remains too short to establish definitive trends.

\paragraph{Effect on Seafood Production}\label{effect-on-seafood-production}

With the poor or unknown stock status of many managed species, the decline in commercial seafood landings in the Gulf of Maine most likely reflects lower catch quotas implemented to rebuild overfished stocks, as well as market dynamics.

The recent decline in \href{Fig.\%20/ref\%7Bfig:rec-landings\%7D}{recreational seafood harvest} is likely associated with a combination of targeted management actions, shifting social behaviors, data collection changes, and potentially low biomass of targeted species. The decline in recreational shark landings can be attributed to management actions intended to reduce fishing mortality on mako sharks. The lower than average landings since 2018 for species other than sharks could be driven by either changes in fishing behavior or a change in NOAA Fisheries' Marine Recreational Information Program survey methodology in 2018. The decline in recreational seafood harvest may also be linked to decreases in sustenance fishing.

Future commercial and recreational landings are likely to be driven by environmental changes that require continued monitoring. Fisheries and communities rely on different combinations of stocks,and individual stocks will respond differently to these drivers. Some key drivers include :

\begin{itemize}
\tightlist
\item
  \textbf{Unprecedented Climate Shifts}: Global ocean temperatures have reached record highs (see \hyperref[highlights]{2025 Highlights section}) and the Northeast US shelf has experienced long-term warming.
\item
  \textbf{Distribution Shifts}: Stocks are shifting towards the northeast and into deeper waters throughout the Northeast US Large Marine Ecosystem (see \hyperref[climate-and-ecosystem-change]{Climate Risks section}).
\item
  \textbf{Ecosystem production}: Changes in ecosystem composition and biological production are impacting the stability of the \hyperref[stability]{ecosystem} and pose \hyperref[risks-to-setting-catch-limits]{risks to setting catch limits}.
\item
  \textbf{Community Risks}: Changes in the ecosystem can affect the \hyperref[stability]{stability} of fisheries and pose \hyperref[community-social-and-climate-vulnerability]{risks to fishing communities}.
\end{itemize}

\begin{figure}

{\centering \includegraphics[width=1\linewidth]{../../../AppData/Local/Temp/RtmpIjNXEs/file26447ab95ec} 

}

\caption{Spring (left) and fall (right) surveyed biomass on Georges Bank, with long-term increasing (orange) and decreasing (purple) trends. The shaded area around each annual mean represents 2 standard deviations from the mean.}\label{fig:nefsc-biomass-gb}
\end{figure}

\begin{figure}

{\centering \includegraphics[width=1\linewidth]{../../../AppData/Local/Temp/RtmpIjNXEs/file26443cae72c9} 

}

\caption{Spring (left) and fall (right) surveyed biomass in the Gulf of Maine, with increasing long-term trends (orange). The shaded area around each annual mean represents 2 standard deviations from the mean.}\label{fig:nefsc-biomass-gom}
\end{figure}

\begin{figure}

{\centering \includegraphics[width=1\linewidth]{../../../AppData/Local/Temp/RtmpIjNXEs/file264478a52e8b} 

}

\caption{Spring (left) and fall (right) surveyed biomass from the state of Massachusetts inshore survey, with increasing (orange) and decreasing (purple) long-term trends. The shaded area around each annual mean represents 2 standard deviations from the mean.}\label{fig:mass-biomass}
\end{figure}

\newpage

\subsection{Commercial Profits}\label{commercial-profits}

\subsubsection{Indicators: revenue (a proxy for profits)}\label{indicators-revenue-a-proxy-for-profits}

Total \href{https://noaa-edab.github.io/catalog/comdat.html}{commercial revenues} from all species is below the long-term mean for GB and near the long-term mean for the GOM in 2024 (Fig. \ref{fig:comm-revenue-ne}). In addition, revenue from NEFMC managed species shows a long-term decline in the GOM. GB continues to exhibit a cyclical nature with regards to revenue, largely driven by rotational management of Atlantic sea scallops.

\begin{figure}

{\centering \includegraphics[width=1\linewidth]{../../../AppData/Local/Temp/RtmpIjNXEs/file26446846e66} 

}

\caption{Commercial revenue through 2023 for Georges Bank (top) and the Gulf of Maine (bottom): total (black) and from NEFMC managed species (red), with significant long-term (light purple) and short-term (dark purple) declines. Dashed lines represent the long-term annual
mean.}\label{fig:comm-revenue-ne}
\end{figure}

Revenue earned by harvesting resources is a function of both the quantity landed of each species and the prices paid for landings. Therefore, total revenue patterns can be driven by harvest levels, the mix of species landed, price changes, or a combination of these. The \href{https://noaa-edab.github.io/catalog/bennet.html}{Bennet Indicator} (BI) decomposes revenue change into two parts, one driven by changing quantities (volumes), and a second driven by changing prices. All changes are in relation to a base year (1982). The 1982 base year was selected because that is the first year the relevant data is available and it allows for an extended period of time to evaluate market trends and dynamics.

In the GB region, revenues have been consistently lower than the 1982 baseline throughout the time series. The changes in total revenue in GB was primarily driven by volumes prior to 2010 rather than by prices (Fig.\ref{fig:bennet-ne}).In more recent years, prices have played a larger role in revenue upticks (such as in 2020), but the overall lower than baseline landings have caused a decreasing revenue in the past three years.

In the GOM, revenues have been above the 1982 baseline in all but three years, with the increase being driven more by relatively higher prices rather than landings. Breaking down the GB revenue by guild (Fig. \ref{fig:bennet-all-ne}), both the volume and price trend have been largely driven by benthivores (lobster) and benthos (scallops, quahogs and surfclams). In the GOM region, increased prices for benthivores (lobster) drove the year-over-year increases in overall prices. Benthivores also had a large influence on the overall volume indicator in the GOM.

\begin{figure}

{\centering \includegraphics[width=1\linewidth]{../../../AppData/Local/Temp/RtmpIjNXEs/file264450db3d41} 

}

\caption{Revenue change from 1982 values in 2023 dollars (black); Price (blue), and Volume Indicators (green) for total
commercial landings in Georges Bank (GB: top) and the Gulf of Maine (GOM: bottom).}\label{fig:bennet-ne}
\end{figure}

\begin{figure}

{\centering \includegraphics[width=1\linewidth]{../../../AppData/Local/Temp/RtmpIjNXEs/file264476657487} 

}

\caption{Total price and volume indicators in 2023 dollars (black) for commercial landings, and individual guild contributions to each indicator from Georges Bank (GB: top panels) and the Gulf of Maine (GOM: bottom panels).}\label{fig:bennet-all-ne}
\end{figure}

\begin{figure}

{\centering \includegraphics[width=1\linewidth]{../../../AppData/Local/Temp/RtmpIjNXEs/file26441d5d7b0f} 

}

\caption{Profitability indices for GB (top) and GOM (bottom): cost index (red), profit index (green), and revenue index (blue). Dashed lines represent the long-term annual means for each index. Long-term increasing trend (orange) in the GOM associated with revenue index.}\label{fig:comdat-profit-ne}
\end{figure}

This year, we present new indicators of \href{https://noaa-edab.github.io/catalog/comdat_profit.html}{profitability}: indices of cost, revenue, and profit. In this index, costs pertain to trip costs, excluding labor, estimated for all federal trips in the region. The profit indicator is net-revenue, determined as the difference between trip revenue and trip costs. Trips were spatially allocated to compile regional indices. Indices are presented as values relative to those from 2000, the first year in the dataset.

In the GOM, the profit index closely follows the same trends as the revenue index with the exception of 2010 - 2013 where low costs created a surge in the profit index. In 2024, the GOM profit index returned to near the long-term average with average costs and revenue. For trips in GB, high costs and low revenue had caused a low profits over the last 3 years, but recent drops in costs have helped compensate for low revenue. GB profits have no long-term trend, but a cyclical revenue driven by rotational scallop management can impact profitability.

For New England ports, \href{https://noaa-edab.github.io/catalog/community_climate_vulnerability.html}{total vulnerability} of revenue was moderate in \hyperref[social-and-community-risks]{2024} with no long-term trend (Fig. \ref{fig:climatevul-rev-ne}). This suggests that while New England commercial fishing is moderately reliant on climate-sensitive species, this proportion has not significantly changed since 2000.

\begin{figure}

{\centering \includegraphics[width=1\linewidth]{../../../AppData/Local/Temp/RtmpIjNXEs/file264466df2ea} 

}

\caption{New England region total climate vulnerability of commercial revenue (sum of New England port revenue weighted by species climate vulnerability from Hare et al. 2016). Horizontal colored bars show different climate risk levels.}\label{fig:climatevul-rev-ne}
\end{figure}

\subsubsection{Implications}\label{implications-1}

The overall volume of lobster and scallops, quahogs and surfclams dictates the revenue trends within the GB region. In the GOM, lobster prices are primarily responsible for relatively high revenues over the time series. Notably, both lobsters and scallops are sensitive to ocean warming and acidification and it is important to monitor the effects of these and other ecosystem drivers.

\newpage

\subsection{Recreational Opportunities}\label{recreational-opportunities}

\subsubsection{Indicators: Angler trips, fleet diversity}\label{indicators-angler-trips-fleet-diversity}

\href{https://noaa-edab.github.io/catalog/recdat.html}{Recreational effort} (angler trips) increased from 1982 to 2010, but has since declined to just below the long-term average (Fig. \ref{fig:rec-op-ne}). Recreational fleets are defined as private vessels, shore-based fishing, or party-charter vessels. Recreational fleet diversity, or the relative importance of each fleet type, has remained relatively stable over the latter half of the time series (Fig. \ref{fig:rec-div-ne}). Billfish landings were notably high in 2025 (See \href{https://noaa-edab.github.io/catalog/observation_synthesis_2025.html}{2025 Highlights} Section), but long-term time series are in development.

\begin{figure}

{\centering \includegraphics[width=1\linewidth]{../../../AppData/Local/Temp/RtmpIjNXEs/file2644234060cd} 

}

\caption{Recreational effort (total number of recreational angler trips from 1980-2023) in New England. Derived from MRIP's Effort Time Series Query.}\label{fig:rec-op-ne}
\end{figure}

\begin{figure}

{\centering \includegraphics[width=1\linewidth]{../../../AppData/Local/Temp/RtmpIjNXEs/file2644737134ba} 

}

\caption{Recreational fleet effort diversity from 1980-2023  in New England.}\label{fig:rec-div-ne}
\end{figure}

\subsubsection{Implications}\label{implications-2}

The absence of a long term trend in recreational angler trips and fleet effort diversity suggests relative stability in the overall number of recreational opportunities in the region.

\newpage

\subsection{Social and Community Risks}\label{social-and-community-risks}

Fisheries management seeks to provide for sustained participation of fishing communities and to avoid adverse economic impacts to fishing communities. A new \href{https://noaa-edab.github.io/catalog/engagement.html?q=engage\#introduction-to-indicator-72}{composite indicator} (Port Commercial Fishing Activity Indicator or PCFA) utilizes NOAA data on dealers, fish landings, and commercial permits to explore trends in commercial fishing activity over time in top ports. This information can be used to understand how changes in fish stocks, regulations, and other social-ecological factors may have disparately impacted ports throughout the Greater Atlantic region.

The recreational engagement index has not been updated from last year and will be updated with similar methods as PCFAI in future reports. The recreational \href{https://noaa-edab.github.io/catalog/engagement.html}{engagement} index demonstrates participation levels in recreational fishing in a given community relative to other coastal communities in a region.

The Community Social Vulnerability Indicators (CSVI) utilize U.S. Census American Community Survey data to describe social characteristics at the municipality level (i.e., not just the fishing community) and provide context for the municipalities utilized by commercial fishing industry participants. Fishing industry participants that live in and/or utilize resources in municipalities with relatively concerning socio-demographic conditions may be more vulnerable to changes. The personal disruption index addresses factors that reduce adaptability to change such as unemployment or educational level. The poverty index is a composite index that indicates a community's financial standing relative to other communities. The population composition index characterizes groups within communities that may be more vulnerable to change. CSVI information for communities highlighted in the PCFA and recreational engagement index have been updated with the most recent census data.

Coastal fishing communities worldwide have or are likely to experience social, economic, and cultural impacts from climate change, both negative (e.g., loss of infrastructure, fish stock decline) and positive (e.g., increased abundance of valuable species). Changes in marine fisheries as a consequence of climate change will require adaptation by coastal fishing communities and fisheries managers alike. The Community Environmental Variability Risk Indicators (CEVRI) were developed to help examine trends in risk related to dependence on species vulnerable to climate and environmental changes.

\subsubsection{Indicators: Port Commercial Fishing Activity and Community Social Vulnerability}\label{indicators-port-commercial-fishing-activity-and-community-social-vulnerability}

The \href{https://noaa-edab.github.io/catalog/engagement.html?q=engage\#introduction-to-indicator-72}{Port Commercial Fishing Activity Indicator (PCFA)} highlights significant shifts in industry engagement across major regional ports. New Bedford, Gloucester, and Boston, MA and Portland, ME, have lower fishing activity compared to their 2007--2011 averages. Because New Bedford and Boston also rank medium-high for socio-demographic vulnerability, industry participants in these municipalities face a higher risk from these changing conditions. Conversely, several communities show substantial growth in fishing activity. Chatham, MA, along with Stonington, Friendship, Harpswell, ME, are seeing increased port activity since 2007-2011.

Of the top 10 most active recreational communities, only Seabrook, NH had medium or higher ranks for at least one socio-demographic indicator (Table \ref{tab:recvultab}) examined here (poverty, personal disruption, population composition). This suggests that future changes to recreational fishing conditions may disproportionately impact Seabrook.

\begin{figure}

{\centering \includegraphics[width=1\linewidth]{../../../AppData/Local/Temp/RtmpIjNXEs/file26443ddbaf8} 

}

\caption{Port Commercial Fishing Activity Indicator scores over time for top commercially active fishing ports in New England.}\label{fig:commercial-engagement-ne}
\end{figure}

\global\setlength{\Oldarrayrulewidth}{\arrayrulewidth}

\global\setlength{\Oldtabcolsep}{\tabcolsep}

\setlength{\tabcolsep}{2pt}

\renewcommand*{\arraystretch}{1.5}



\providecommand{\ascline}[3]{\noalign{\global\arrayrulewidth #1}\arrayrulecolor[HTML]{#2}\cline{#3}}

\begin{longtable}[c]{|p{1.89in}|p{1.38in}|p{1.61in}|p{0.80in}}

\caption{Community\ Vulnerability\ factors\ for\ the\ top\ 10\ commercially\ active\ New\ England\ Ports}\label{tab:comvultab}\\

\ascline{1.5pt}{666666}{1-4}

\multicolumn{1}{>{\raggedright}m{\dimexpr 1.89in+0\tabcolsep}}{\textcolor[HTML]{000000}{\fontsize{9}{9}\selectfont{Community}}} & \multicolumn{1}{>{\raggedright}m{\dimexpr 1.38in+0\tabcolsep}}{\textcolor[HTML]{000000}{\fontsize{9}{9}\selectfont{Personal\ Disruption}}} & \multicolumn{1}{>{\raggedright}m{\dimexpr 1.61in+0\tabcolsep}}{\textcolor[HTML]{000000}{\fontsize{9}{9}\selectfont{Population\ Composition}}} & \multicolumn{1}{>{\raggedright}m{\dimexpr 0.8in+0\tabcolsep}}{\textcolor[HTML]{000000}{\fontsize{9}{9}\selectfont{Poverty}}} \\

\ascline{1.5pt}{666666}{1-4}\endfirsthead \caption[]{Community\ Vulnerability\ factors\ for\ the\ top\ 10\ commercially\ active\ New\ England\ Ports}\label{tab:comvultab}\\

\ascline{1.5pt}{666666}{1-4}

\multicolumn{1}{>{\raggedright}m{\dimexpr 1.89in+0\tabcolsep}}{\textcolor[HTML]{000000}{\fontsize{9}{9}\selectfont{Community}}} & \multicolumn{1}{>{\raggedright}m{\dimexpr 1.38in+0\tabcolsep}}{\textcolor[HTML]{000000}{\fontsize{9}{9}\selectfont{Personal\ Disruption}}} & \multicolumn{1}{>{\raggedright}m{\dimexpr 1.61in+0\tabcolsep}}{\textcolor[HTML]{000000}{\fontsize{9}{9}\selectfont{Population\ Composition}}} & \multicolumn{1}{>{\raggedright}m{\dimexpr 0.8in+0\tabcolsep}}{\textcolor[HTML]{000000}{\fontsize{9}{9}\selectfont{Poverty}}} \\

\ascline{1.5pt}{666666}{1-4}\endhead



\multicolumn{1}{>{\raggedright}m{\dimexpr 1.89in+0\tabcolsep}}{\textcolor[HTML]{000000}{\fontsize{9}{9}\selectfont{New\ Bedford,\ MA}}} & \multicolumn{1}{>{\raggedright}m{\dimexpr 1.38in+0\tabcolsep}}{\textcolor[HTML]{000000}{\fontsize{9}{9}\selectfont{high}}} & \multicolumn{1}{>{\raggedright}m{\dimexpr 1.61in+0\tabcolsep}}{\textcolor[HTML]{000000}{\fontsize{9}{9}\selectfont{med\ high}}} & \multicolumn{1}{>{\raggedright}m{\dimexpr 0.8in+0\tabcolsep}}{\textcolor[HTML]{000000}{\fontsize{9}{9}\selectfont{med\ high}}} \\





\multicolumn{1}{>{\raggedright}m{\dimexpr 1.89in+0\tabcolsep}}{\textcolor[HTML]{000000}{\fontsize{9}{9}\selectfont{Boston,\ MA}}} & \multicolumn{1}{>{\raggedright}m{\dimexpr 1.38in+0\tabcolsep}}{\textcolor[HTML]{000000}{\fontsize{9}{9}\selectfont{med}}} & \multicolumn{1}{>{\raggedright}m{\dimexpr 1.61in+0\tabcolsep}}{\textcolor[HTML]{000000}{\fontsize{9}{9}\selectfont{med\ high}}} & \multicolumn{1}{>{\raggedright}m{\dimexpr 0.8in+0\tabcolsep}}{\textcolor[HTML]{000000}{\fontsize{9}{9}\selectfont{med\ high}}} \\





\multicolumn{1}{>{\raggedright}m{\dimexpr 1.89in+0\tabcolsep}}{\textcolor[HTML]{000000}{\fontsize{9}{9}\selectfont{Gloucester,\ MA}}} & \multicolumn{1}{>{\raggedright}m{\dimexpr 1.38in+0\tabcolsep}}{\textcolor[HTML]{000000}{\fontsize{9}{9}\selectfont{low}}} & \multicolumn{1}{>{\raggedright}m{\dimexpr 1.61in+0\tabcolsep}}{\textcolor[HTML]{000000}{\fontsize{9}{9}\selectfont{low}}} & \multicolumn{1}{>{\raggedright}m{\dimexpr 0.8in+0\tabcolsep}}{\textcolor[HTML]{000000}{\fontsize{9}{9}\selectfont{med}}} \\





\multicolumn{1}{>{\raggedright}m{\dimexpr 1.89in+0\tabcolsep}}{\textcolor[HTML]{000000}{\fontsize{9}{9}\selectfont{Rockland,\ ME}}} & \multicolumn{1}{>{\raggedright}m{\dimexpr 1.38in+0\tabcolsep}}{\textcolor[HTML]{000000}{\fontsize{9}{9}\selectfont{med}}} & \multicolumn{1}{>{\raggedright}m{\dimexpr 1.61in+0\tabcolsep}}{\textcolor[HTML]{000000}{\fontsize{9}{9}\selectfont{low}}} & \multicolumn{1}{>{\raggedright}m{\dimexpr 0.8in+0\tabcolsep}}{\textcolor[HTML]{000000}{\fontsize{9}{9}\selectfont{low}}} \\





\multicolumn{1}{>{\raggedright}m{\dimexpr 1.89in+0\tabcolsep}}{\textcolor[HTML]{000000}{\fontsize{9}{9}\selectfont{Jonesport,\ ME}}} & \multicolumn{1}{>{\raggedright}m{\dimexpr 1.38in+0\tabcolsep}}{\textcolor[HTML]{000000}{\fontsize{9}{9}\selectfont{med}}} & \multicolumn{1}{>{\raggedright}m{\dimexpr 1.61in+0\tabcolsep}}{\textcolor[HTML]{000000}{\fontsize{9}{9}\selectfont{low}}} & \multicolumn{1}{>{\raggedright}m{\dimexpr 0.8in+0\tabcolsep}}{\textcolor[HTML]{000000}{\fontsize{9}{9}\selectfont{low}}} \\





\multicolumn{1}{>{\raggedright}m{\dimexpr 1.89in+0\tabcolsep}}{\textcolor[HTML]{000000}{\fontsize{9}{9}\selectfont{Chatham,\ MA}}} & \multicolumn{1}{>{\raggedright}m{\dimexpr 1.38in+0\tabcolsep}}{\textcolor[HTML]{000000}{\fontsize{9}{9}\selectfont{low}}} & \multicolumn{1}{>{\raggedright}m{\dimexpr 1.61in+0\tabcolsep}}{\textcolor[HTML]{000000}{\fontsize{9}{9}\selectfont{low}}} & \multicolumn{1}{>{\raggedright}m{\dimexpr 0.8in+0\tabcolsep}}{\textcolor[HTML]{000000}{\fontsize{9}{9}\selectfont{low}}} \\





\multicolumn{1}{>{\raggedright}m{\dimexpr 1.89in+0\tabcolsep}}{\textcolor[HTML]{000000}{\fontsize{9}{9}\selectfont{Portland,\ ME}}} & \multicolumn{1}{>{\raggedright}m{\dimexpr 1.38in+0\tabcolsep}}{\textcolor[HTML]{000000}{\fontsize{9}{9}\selectfont{low}}} & \multicolumn{1}{>{\raggedright}m{\dimexpr 1.61in+0\tabcolsep}}{\textcolor[HTML]{000000}{\fontsize{9}{9}\selectfont{low}}} & \multicolumn{1}{>{\raggedright}m{\dimexpr 0.8in+0\tabcolsep}}{\textcolor[HTML]{000000}{\fontsize{9}{9}\selectfont{low}}} \\





\multicolumn{1}{>{\raggedright}m{\dimexpr 1.89in+0\tabcolsep}}{\textcolor[HTML]{000000}{\fontsize{9}{9}\selectfont{Stonington,\ ME}}} & \multicolumn{1}{>{\raggedright}m{\dimexpr 1.38in+0\tabcolsep}}{\textcolor[HTML]{000000}{\fontsize{9}{9}\selectfont{low}}} & \multicolumn{1}{>{\raggedright}m{\dimexpr 1.61in+0\tabcolsep}}{\textcolor[HTML]{000000}{\fontsize{9}{9}\selectfont{low}}} & \multicolumn{1}{>{\raggedright}m{\dimexpr 0.8in+0\tabcolsep}}{\textcolor[HTML]{000000}{\fontsize{9}{9}\selectfont{low}}} \\





\multicolumn{1}{>{\raggedright}m{\dimexpr 1.89in+0\tabcolsep}}{\textcolor[HTML]{000000}{\fontsize{9}{9}\selectfont{Friendship,\ ME}}} & \multicolumn{1}{>{\raggedright}m{\dimexpr 1.38in+0\tabcolsep}}{\textcolor[HTML]{000000}{\fontsize{9}{9}\selectfont{low}}} & \multicolumn{1}{>{\raggedright}m{\dimexpr 1.61in+0\tabcolsep}}{\textcolor[HTML]{000000}{\fontsize{9}{9}\selectfont{low}}} & \multicolumn{1}{>{\raggedright}m{\dimexpr 0.8in+0\tabcolsep}}{\textcolor[HTML]{000000}{\fontsize{9}{9}\selectfont{low}}} \\





\multicolumn{1}{>{\raggedright}m{\dimexpr 1.89in+0\tabcolsep}}{\textcolor[HTML]{000000}{\fontsize{9}{9}\selectfont{Narragansett/Point\ Judith,\ RI}}} & \multicolumn{1}{>{\raggedright}m{\dimexpr 1.38in+0\tabcolsep}}{\textcolor[HTML]{000000}{\fontsize{9}{9}\selectfont{low}}} & \multicolumn{1}{>{\raggedright}m{\dimexpr 1.61in+0\tabcolsep}}{\textcolor[HTML]{000000}{\fontsize{9}{9}\selectfont{low}}} & \multicolumn{1}{>{\raggedright}m{\dimexpr 0.8in+0\tabcolsep}}{\textcolor[HTML]{000000}{\fontsize{9}{9}\selectfont{low}}} \\

\ascline{1.5pt}{666666}{1-4}



\end{longtable}



\arrayrulecolor[HTML]{000000}

\global\setlength{\arrayrulewidth}{\Oldarrayrulewidth}

\global\setlength{\tabcolsep}{\Oldtabcolsep}

\renewcommand*{\arraystretch}{1}

\begin{figure}

{\centering \includegraphics[width=1\linewidth]{../../../AppData/Local/Temp/RtmpIjNXEs/file264439565e3d} 

}

\caption{Recreational engagement and population relative engagement with labels for the top recreationally engaged fishing communities in New England (last updated 2024).}\label{fig:recreational-engagement-ne}
\end{figure}

\global\setlength{\Oldarrayrulewidth}{\arrayrulewidth}

\global\setlength{\Oldtabcolsep}{\tabcolsep}

\setlength{\tabcolsep}{2pt}

\renewcommand*{\arraystretch}{1.5}



\providecommand{\ascline}[3]{\noalign{\global\arrayrulewidth #1}\arrayrulecolor[HTML]{#2}\cline{#3}}

\begin{longtable}[c]{|p{2.57in}|p{1.38in}|p{1.61in}|p{0.71in}}

\caption{Community\ Vulnerability\ factors\ for\ the\ top\ 10\ recreationally\ engaged\ New\ England\ Ports}\label{tab:recvultab}\\

\ascline{1.5pt}{666666}{1-4}

\multicolumn{1}{>{\raggedright}m{\dimexpr 2.57in+0\tabcolsep}}{\textcolor[HTML]{000000}{\fontsize{9}{9}\selectfont{Community}}} & \multicolumn{1}{>{\raggedright}m{\dimexpr 1.38in+0\tabcolsep}}{\textcolor[HTML]{000000}{\fontsize{9}{9}\selectfont{Personal\ Disruption}}} & \multicolumn{1}{>{\raggedright}m{\dimexpr 1.61in+0\tabcolsep}}{\textcolor[HTML]{000000}{\fontsize{9}{9}\selectfont{Population\ Composition}}} & \multicolumn{1}{>{\raggedright}m{\dimexpr 0.71in+0\tabcolsep}}{\textcolor[HTML]{000000}{\fontsize{9}{9}\selectfont{Poverty}}} \\

\ascline{1.5pt}{666666}{1-4}\endfirsthead \caption[]{Community\ Vulnerability\ factors\ for\ the\ top\ 10\ recreationally\ engaged\ New\ England\ Ports}\label{tab:recvultab}\\

\ascline{1.5pt}{666666}{1-4}

\multicolumn{1}{>{\raggedright}m{\dimexpr 2.57in+0\tabcolsep}}{\textcolor[HTML]{000000}{\fontsize{9}{9}\selectfont{Community}}} & \multicolumn{1}{>{\raggedright}m{\dimexpr 1.38in+0\tabcolsep}}{\textcolor[HTML]{000000}{\fontsize{9}{9}\selectfont{Personal\ Disruption}}} & \multicolumn{1}{>{\raggedright}m{\dimexpr 1.61in+0\tabcolsep}}{\textcolor[HTML]{000000}{\fontsize{9}{9}\selectfont{Population\ Composition}}} & \multicolumn{1}{>{\raggedright}m{\dimexpr 0.71in+0\tabcolsep}}{\textcolor[HTML]{000000}{\fontsize{9}{9}\selectfont{Poverty}}} \\

\ascline{1.5pt}{666666}{1-4}\endhead



\multicolumn{1}{>{\raggedright}m{\dimexpr 2.57in+0\tabcolsep}}{\textcolor[HTML]{000000}{\fontsize{9}{9}\selectfont{Seabrook,\ NH}}} & \multicolumn{1}{>{\raggedright}m{\dimexpr 1.38in+0\tabcolsep}}{\textcolor[HTML]{000000}{\fontsize{9}{9}\selectfont{med}}} & \multicolumn{1}{>{\raggedright}m{\dimexpr 1.61in+0\tabcolsep}}{\textcolor[HTML]{000000}{\fontsize{9}{9}\selectfont{low}}} & \multicolumn{1}{>{\raggedright}m{\dimexpr 0.71in+0\tabcolsep}}{\textcolor[HTML]{000000}{\fontsize{9}{9}\selectfont{med}}} \\





\multicolumn{1}{>{\raggedright}m{\dimexpr 2.57in+0\tabcolsep}}{\textcolor[HTML]{000000}{\fontsize{9}{9}\selectfont{Old\ Saybrook,\ CT}}} & \multicolumn{1}{>{\raggedright}m{\dimexpr 1.38in+0\tabcolsep}}{\textcolor[HTML]{000000}{\fontsize{9}{9}\selectfont{low}}} & \multicolumn{1}{>{\raggedright}m{\dimexpr 1.61in+0\tabcolsep}}{\textcolor[HTML]{000000}{\fontsize{9}{9}\selectfont{low}}} & \multicolumn{1}{>{\raggedright}m{\dimexpr 0.71in+0\tabcolsep}}{\textcolor[HTML]{000000}{\fontsize{9}{9}\selectfont{low}}} \\





\multicolumn{1}{>{\raggedright}m{\dimexpr 2.57in+0\tabcolsep}}{\textcolor[HTML]{000000}{\fontsize{9}{9}\selectfont{Waterford,\ CT}}} & \multicolumn{1}{>{\raggedright}m{\dimexpr 1.38in+0\tabcolsep}}{\textcolor[HTML]{000000}{\fontsize{9}{9}\selectfont{low}}} & \multicolumn{1}{>{\raggedright}m{\dimexpr 1.61in+0\tabcolsep}}{\textcolor[HTML]{000000}{\fontsize{9}{9}\selectfont{low}}} & \multicolumn{1}{>{\raggedright}m{\dimexpr 0.71in+0\tabcolsep}}{\textcolor[HTML]{000000}{\fontsize{9}{9}\selectfont{low}}} \\





\multicolumn{1}{>{\raggedright}m{\dimexpr 2.57in+0\tabcolsep}}{\textcolor[HTML]{000000}{\fontsize{9}{9}\selectfont{Harwich\ Port,\ MA}}} & \multicolumn{1}{>{\raggedright}m{\dimexpr 1.38in+0\tabcolsep}}{\textcolor[HTML]{000000}{\fontsize{9}{9}\selectfont{low}}} & \multicolumn{1}{>{\raggedright}m{\dimexpr 1.61in+0\tabcolsep}}{\textcolor[HTML]{000000}{\fontsize{9}{9}\selectfont{low}}} & \multicolumn{1}{>{\raggedright}m{\dimexpr 0.71in+0\tabcolsep}}{\textcolor[HTML]{000000}{\fontsize{9}{9}\selectfont{low}}} \\





\multicolumn{1}{>{\raggedright}m{\dimexpr 2.57in+0\tabcolsep}}{\textcolor[HTML]{000000}{\fontsize{9}{9}\selectfont{Newburyport,\ MA}}} & \multicolumn{1}{>{\raggedright}m{\dimexpr 1.38in+0\tabcolsep}}{\textcolor[HTML]{000000}{\fontsize{9}{9}\selectfont{low}}} & \multicolumn{1}{>{\raggedright}m{\dimexpr 1.61in+0\tabcolsep}}{\textcolor[HTML]{000000}{\fontsize{9}{9}\selectfont{low}}} & \multicolumn{1}{>{\raggedright}m{\dimexpr 0.71in+0\tabcolsep}}{\textcolor[HTML]{000000}{\fontsize{9}{9}\selectfont{low}}} \\





\multicolumn{1}{>{\raggedright}m{\dimexpr 2.57in+0\tabcolsep}}{\textcolor[HTML]{000000}{\fontsize{9}{9}\selectfont{Barnstable\ Town,\ MA}}} & \multicolumn{1}{>{\raggedright}m{\dimexpr 1.38in+0\tabcolsep}}{\textcolor[HTML]{000000}{\fontsize{9}{9}\selectfont{low}}} & \multicolumn{1}{>{\raggedright}m{\dimexpr 1.61in+0\tabcolsep}}{\textcolor[HTML]{000000}{\fontsize{9}{9}\selectfont{low}}} & \multicolumn{1}{>{\raggedright}m{\dimexpr 0.71in+0\tabcolsep}}{\textcolor[HTML]{000000}{\fontsize{9}{9}\selectfont{low}}} \\





\multicolumn{1}{>{\raggedright}m{\dimexpr 2.57in+0\tabcolsep}}{\textcolor[HTML]{000000}{\fontsize{9}{9}\selectfont{Sandwich/East\ Sandwich/Forestdale,\ MA}}} & \multicolumn{1}{>{\raggedright}m{\dimexpr 1.38in+0\tabcolsep}}{\textcolor[HTML]{000000}{\fontsize{9}{9}\selectfont{low}}} & \multicolumn{1}{>{\raggedright}m{\dimexpr 1.61in+0\tabcolsep}}{\textcolor[HTML]{000000}{\fontsize{9}{9}\selectfont{low}}} & \multicolumn{1}{>{\raggedright}m{\dimexpr 0.71in+0\tabcolsep}}{\textcolor[HTML]{000000}{\fontsize{9}{9}\selectfont{low}}} \\





\multicolumn{1}{>{\raggedright}m{\dimexpr 2.57in+0\tabcolsep}}{\textcolor[HTML]{000000}{\fontsize{9}{9}\selectfont{Westport,\ MA}}} & \multicolumn{1}{>{\raggedright}m{\dimexpr 1.38in+0\tabcolsep}}{\textcolor[HTML]{000000}{\fontsize{9}{9}\selectfont{low}}} & \multicolumn{1}{>{\raggedright}m{\dimexpr 1.61in+0\tabcolsep}}{\textcolor[HTML]{000000}{\fontsize{9}{9}\selectfont{low}}} & \multicolumn{1}{>{\raggedright}m{\dimexpr 0.71in+0\tabcolsep}}{\textcolor[HTML]{000000}{\fontsize{9}{9}\selectfont{low}}} \\





\multicolumn{1}{>{\raggedright}m{\dimexpr 2.57in+0\tabcolsep}}{\textcolor[HTML]{000000}{\fontsize{9}{9}\selectfont{Hampton,\ NH}}} & \multicolumn{1}{>{\raggedright}m{\dimexpr 1.38in+0\tabcolsep}}{\textcolor[HTML]{000000}{\fontsize{9}{9}\selectfont{low}}} & \multicolumn{1}{>{\raggedright}m{\dimexpr 1.61in+0\tabcolsep}}{\textcolor[HTML]{000000}{\fontsize{9}{9}\selectfont{low}}} & \multicolumn{1}{>{\raggedright}m{\dimexpr 0.71in+0\tabcolsep}}{\textcolor[HTML]{000000}{\fontsize{9}{9}\selectfont{low}}} \\





\multicolumn{1}{>{\raggedright}m{\dimexpr 2.57in+0\tabcolsep}}{\textcolor[HTML]{000000}{\fontsize{9}{9}\selectfont{Narragansett/Point\ Judith,\ RI}}} & \multicolumn{1}{>{\raggedright}m{\dimexpr 1.38in+0\tabcolsep}}{\textcolor[HTML]{000000}{\fontsize{9}{9}\selectfont{low}}} & \multicolumn{1}{>{\raggedright}m{\dimexpr 1.61in+0\tabcolsep}}{\textcolor[HTML]{000000}{\fontsize{9}{9}\selectfont{low}}} & \multicolumn{1}{>{\raggedright}m{\dimexpr 0.71in+0\tabcolsep}}{\textcolor[HTML]{000000}{\fontsize{9}{9}\selectfont{low}}} \\

\ascline{1.5pt}{666666}{1-4}



\end{longtable}



\arrayrulecolor[HTML]{000000}

\global\setlength{\arrayrulewidth}{\Oldarrayrulewidth}

\global\setlength{\tabcolsep}{\Oldtabcolsep}

\renewcommand*{\arraystretch}{1}

\newpage

\subsubsection{Indicators: Community Environmental Variability Risk in New England}\label{indicators-community-environmental-variability-risk-in-new-england}

\href{https://noaa-edab.github.io/catalog/community_climate_vulnerability.html}{Community Environmental Variability Risk Indicators} (CEVRI) measure risk by linking commercial landings and revenue to specific climate sensitivity factors, including temperature, ocean acidification, and stock status using the Climate Vulnerability Assessment (CVA) scores (based on Hare et al.~2016). These indicators calculate total sensitivity and vulnerability scores based on a community's dependence on species vulnerable to climate change. Risk scores range from low (1) to high (4), increasing as a community relies more heavily on species at higher risk from environmental shifts.While there is no long-term trend in risk across New England communities, the proportion of communities with moderate risk is decreasing and shifting more towards high or very high risk scores (Fig. \ref{fig:commvulex-ne}). This shift demonstrates that regional communities are increasing their dependence on species that are highly vulnerable to changing ocean conditions for their commercial revenue. Strategies for management should account for this increased reliance on climate-sensitive stocks.

\begin{figure}

{\centering \includegraphics[width=1\linewidth]{../../../AppData/Local/Temp/RtmpIjNXEs/file264428111da7} 

}

\caption{Proportion of New England communities at each revenue climate vulnerability level over time. Total climate vulnerability ranges from low (green), moderate (yellow), high (orange), to very high (red).}\label{fig:commvulex-ne}
\end{figure}

\subsubsection{Implications}\label{implications-3}

Social and demographic indicators highlight potential vulnerabilities in New England's most active commercial fishing ports. Industry participants in these locations face increased risk from shifting fishing patterns, whether driven by new regulations or broader ecosystem changes. Because many of these primary communities show medium to high socio-demographic risk, they may lack the necessary resources to adapt effectively to industry transitions

\newpage

\subsection{Protected Species}\label{protected-species}

Fishery management objectives for protected species generally focus on reducing threats and on habitat conservation/restoration. Specific actions include managing bycatch to remain below potential biological removal (PBR) thresholds, recovering endangered populations, and monitoring unusual mortality events (UMEs). Protected species include marine mammals protected under the Marine Mammal Protection Act, endangered and threatened species protected under the Endangered Species Act, and migratory birds protected under the Migratory Bird Treaty Act. In the Northeast U.S., endangered/threatened species include Atlantic salmon, Atlantic and shortnose sturgeon, all sea turtle species, giant manta ray, oceanic whitetip shark, and five baleen whales. Here we report on performance relative to these objectives, as well as how observed and predicted ecosystem changes in the Northeast U.S may impact these objectives in the future.

\subsubsection{Indicators: bycatch, population (adult and juvenile) numbers, mortalities}\label{indicators-bycatch-population-adult-and-juvenile-numbers-mortalities}

The management objective for \href{https://noaa-edab.github.io/catalog/harborporpoise.html}{harbor porpoise} has been met, as the average index (Fig. \ref{fig:harborporpoise}) remains below the current PBR threshold.

\begin{figure}

{\centering \includegraphics[width=1\linewidth]{../../../AppData/Local/Temp/RtmpIjNXEs/file2644334f300a} 

}

\caption{Harbor porpoise average bycatch estimate for Mid-Atlantic and New England gillnet fisheries (blue, confidence interval shaded) and the potential biological removal (red). The dashed line (black) represents the annual estimated bycatch. }\label{fig:harborporpoise}
\end{figure}

The annual estimate for gray seal bycatch, most of which occurs in New England, has generally declined since 2019, in part driven by declining gillnet landings. Although, post-2019 estimates have greater uncertainty stemming from low observer coverage. The U.S. and Canadian range-wide PBR for \href{https://noaa-edab.github.io/catalog/grayseal.html}{gray seals} is 12,052. Despite the PBR for the portion of this stock in U.S. waters being reduced to 756 animals bycatch (Fig. \ref{fig:grayseal}), but due to incomplete data on anthropogenic mortality and serious injury, bycatch for this stock is still considered unlikely to exceed the range-wide PBR. Thus the bycatch management objective gray seals has been met.

\begin{figure}

{\centering \includegraphics[width=1\linewidth]{../../../AppData/Local/Temp/RtmpIjNXEs/file26443a5bc6b} 

}

\caption{Gray seal five-year average bycatch estimate for New England and Mid-Atlantic U.S. gillnet fisheries (blue, with confidence interval shaded) and the potential U.S. biological removal (red). The range-wide PBR, including both U.S. and Canadian portions of the population, is 12,052 in the draft 2024 SAR. The dashed line (black) represents the annual estimated bycatch.}\label{fig:grayseal}
\end{figure}

The \href{https://noaa-edab.github.io/catalog/narw.html}{North Atlantic right whale population} was on a recovery trajectory until 2010, but has since declined (Fig. \ref{fig:narw-abundance}), slowing since 2020. The right whale population continues to experience annual mortalities above recovery thresholds. Reduced survival rates of adult females lead to diverging abundance trends between sexes. It is estimated that there are approximately 70 adult females remaining in the population.

\begin{figure}

{\centering \includegraphics[width=1\linewidth]{../../../AppData/Local/Temp/RtmpIjNXEs/file26447224af6} 

}

\caption{Estimated North Atlantic right whale abundance on the Northeast Shelf. 95\% confidence interval shaded in gray around the line. Analysis is based on methods by Pace, Corkeron, and Kraus (2017), as documented most recently by Linden (2025).}\label{fig:narw-abundance}
\end{figure}

North Atlantic right whale \href{https://noaa-edab.github.io/catalog/narw.html}{calf counts} have generally declined after 2009 to the point of having zero new calves observed in 2018 (Fig. \ref{fig:NARW-calf-abundance}). However, since 2020, calf births have been closer to the long-term average, with 11 calves born in 2025.

\begin{figure}

{\centering \includegraphics[width=1\linewidth]{../../../AppData/Local/Temp/RtmpIjNXEs/file26444db42d1d} 

}

\caption{Number of North Atlantic right whale calf births since 1980. Calf birth estimates are available in Linden (2025).}\label{fig:NARW-calf-abundance}
\end{figure}

Human interaction from entanglements or vessel strikes remains the primary cause of death in an ongoing North Atlantic right whale Unusual Mortality Event (UME) since 2017. As of 5 January 2026, the UME includes 168 individual whales: 41 confirmed mortalities (19 US; 22 Canada), 40 serious injuries, and 87 sublethal injuries or illnesses. Recent research suggests that many mortalities go unobserved and the true number of mortalities are about three times the count of the observed mortalities.

There is an ongoing UME for humpback whales (2016-present) and Atlantic minke whales (2018-present); suspected causes include human interactions. A UME for Northeast pinnipeds that began in 2018 for infectious disease is non-active pending closure as of February 2026.

\subsubsection{Implications}\label{implications-4}

Bycatch management measures have been implemented to maintain bycatch below PBR thresholds. The downward trend in harbor porpoise bycatch could also be due to a decrease in harbor porpoise abundance in U.S. waters, reducing their overlap with fisheries, and a decrease in gillnet effort. The increasing trend in 5-year average gray seal bycatch may be related to an increase in the gray seal population (\href{https://noaa-edab.github.io/catalog/seal_pups.html}{U.S. pup counts}), supported by the dramatic rise over the last three decades in observed numbers of gray seal pups born at U.S. breeding sites plus an increase in adult seals at the breeding sites, some of which are supplemented by Canadian adults.

\begin{figure}

{\centering \includegraphics[width=1\linewidth]{../../../AppData/Local/Temp/RtmpIjNXEs/file26446dc67f0a} 

}

\caption{Estimated number of gray seal pups born at four United States pupping colonies at various times from 1988 to 2021. Recreated from Wood et al. 2022 (Figure 5).}\label{fig:seals}
\end{figure}

Strong evidence exists to suggest that interactions between right whales and both the fixed gear fisheries in the U.S. and Canada and vessel strikes in the U.S. are contributing substantially to the decline of the species. Further, right whale distribution has changed since 2010. \href{https://noaa-edab.github.io/catalog/narw.html}{Recent research} suggests that recent climate driven changes in ocean circulation have resulted in right whale distribution changes driven by increased warm water influx through the Northeast Channel, which has reduced the \href{https://noaa-edab.github.io/catalog/preyfield_energy.html}{primary right whale prey} (the copepod \emph{Calanus finmarchicus}) in the central and eastern portions of the Gulf of Maine. Additional potential stressors include offshore wind development, which overlaps with important habitat areas used year-round by right whales, including mother and calf migration corridors and foraging habitat. Additional information can be found in the \hyperref[other-ocean-uses-offshore-wind]{offshore wind risks section}.

The UMEs are under investigation and are likely the result of multiple drivers. For all large whale UMEs, human interaction appears to have contributed to increased mortalities, although investigations are not complete.

A climate vulnerability assessment (Lettrich et al.~2023) is published for Atlantic and Gulf marine mammal populations.

\newpage

\subsection{Stability}\label{stability}

\paragraph{Fishery Stability}\label{fishery-stability}

Fisheries in Georges Bank and the Gulf of Maine are dominated by single species. Total landings are declining in both regions, although overall revenue does not have a long-term trend. Revenue from Council-managed fisheries in the GOM and GB have declined over time. However, \href{https://noaa-edab.github.io/catalog/effective_sweptarea.html}{revenue per unit effort} remains steady or increasing over time for most gear types, indicating financial viability of current fishing operations. Although the effective number of species being landed in the commercial fleet rebounded slightly from the historical low of 2021, the \href{https://noaa-edab.github.io/catalog/recdat.html}{diversity in catch} is still well below the series average, indicating increasing reliance on a smaller number of species. Commercial fishery fleet count is also below the time series average due to varying barriers to enter and invest into the fleet. While some opportunity to diversify catch has allowed crew and vessel owners to continue at a sustainable rate, other barriers such as shifting species distribution and population shifts leave commercial crew and vessel owners in vulnerable positions to adapt to these changes. In Georges Bank, cyclic landings and revenue patterns are driven by scallops and decrease the predictability of earnings from year to year. In the Gulf of Maine, landings and revenue are driven by lobster. The increasing importance of lobster over time is mirrored in the decreasing contribution of Council-managed fisheries to Gulf of Maine total landings and revenue.

Results from the \href{https://noaa-edab.github.io/catalog/crew_survey.html}{Crew Survey} suggest many commercial fishing crews in New England are dissatisfied with the predictability of their earnings, the amount of time away from home, and the physical fatigue and personal health impacts from the job. Additionally, the survey results demonstrate evidence of aging or ``graying'' of the fleet in New England, which combined with a lack of new entrants to the industry suggests that participation in commercial fishing is declining across the region.

\href{https://noaa-edab.github.io/catalog/communities_at_sea.html}{Communities at Sea} indicators show a decline in the number of New England fishing communities since 1996, suggesting a consolidation of fishing operations and employment concentrated into fewer ports. Fishing days on trawlers, a proxy for employment, has also declined over this time period, while employment in lobster communities has increased. Adaptive capacity indicators show that the ability to shift target species and fishing grounds varies by community, which is most limited in lobster potting communities. The Communities at Sea indicators combined with the declining fleet count and declining overall landings, suggests a reduced capacity for New England commercial fisheries to adapt to future change.

The number of \href{https://noaa-edab.github.io/catalog/recdat.html}{recreational trips} is below average, although there is no long-term trend, and recreational landings have been declining. Low recreational landings may also be driven by a shift to catch-and-release fishing and stricter shark and large sport fish regulations. Recreational effort diversity is near average with no trend and there has been no shift in angling modes, suggesting steady recreational fishing opportunities. Recreational species catch diversity has increased over time (Fig. \ref{fig:recdat-div-catch-ne}), indicating that anglers are catching a more varied mix of species, likely due to shifting species distributions.

\begin{figure}

{\centering \includegraphics[width=1\linewidth]{../../../AppData/Local/Temp/RtmpIjNXEs/file2644275f7519} 

}

\caption{Species revenue diversity (permit-level species effective Shannon index) in New England.}\label{fig:permit-div}
\end{figure}

\begin{figure}

{\centering \includegraphics[width=1\linewidth]{../../../AppData/Local/Temp/RtmpIjNXEs/file2644168c1e23} 

}

\caption{Diversity of recreational catch in New England, with long term increasing trend (orange). Derived from MRIP's Catch Time Series Query.}\label{fig:recdat-div-catch-ne}
\end{figure}

\paragraph{Ecological Stability}\label{ecological-stability}

Long-term changes in biological production suggest the Georges Bank ecosystem is experiencing a system-wide shift. Total annual \href{https://noaa-edab.github.io/catalog/chl_pp.html}{primary production} (TPP) is a measure of the total amount of carbon (i.e., energy) produced by phytoplankton per year and is variable over time. Zooplankton biomass and the biomass of some groups of fish are also increasing. However, the productivity of managed species has declined over time, suggesting that although the system remains productive, this productivity is driven by non-target species.

\begin{figure}

{\centering \includegraphics[width=1\linewidth]{../../../AppData/Local/Temp/RtmpIjNXEs/file2644af04f9a} 

}

\caption{Total areal annual primary production by ecological production unit (Georges Bank, top; Gulf of Maine, bottom). The dashed line represents the long-term (1998-2023) annual mean.}\label{fig:totpp-ne}
\end{figure}

\newpage

The Gulf of Maine ecosystem has also continued to change over time leading to less predictable ecosystem conditions. Long-term primary productivity has remained relatively constant, but increases in planktivores and euphausiids suggest changing ecosystem dynamics and potential for complex interactions with higher trophic levels. The zooplankton community displays distinct regime shifts in composition corresponding to approximately decadal time scales, with the most recent shift in community composition occurring in 2023. Productivity of managed species has declined, with current levels below average, but it is unclear if that is the result of these ecosystem changes.

\begin{figure}

{\centering \includegraphics[width=1\linewidth]{../../../AppData/Local/Temp/RtmpIjNXEs/file264442e62fc5} 

}

\end{figure}

\begin{figure}

{\centering \includegraphics[width=1\linewidth]{../../../AppData/Local/Temp/RtmpIjNXEs/file26445331550f} 

}

\caption{Adult fish diversity for Georges Bank (top) and in the Gulf of Maine (bottom) with long-term (light orange) and short-term (dark orange) increasing trends, based on expected number of species in a standard number of individuals. Results from survey vessels Albatross (red) and Bigelow (blue) are reported separately due to catchability differences.}\label{fig:exp-n-ne}
\end{figure}

\href{https://noaa-edab.github.io/catalog/finfish_traits.html}{Functional traits}, such as length at maturity, maximum body size, or fecundity, serve to synthesize change in complex, diverse communities by looking beyond species-specific trends. Furthermore, shifts in functional trait distributions for the fish community can indicate changes in ecosystem-scale resilience. There is evidence for shifts in functional trait distributions in New England (Fig. \ref{fig:traits-k-ne}) (Fig. \ref{fig:traits-tl-ne}). George's Bank (GB) displayed few long-term trends other than reductions in fecundity in both fall and spring. The Gulf of Maine (GOM) displayed long-term trends consistent with shifts towards faster life history strategies particularly in the spring finfish community, including younger age and shorter length at maturity, lower fecundity, and faster growth rate. Interestly, the spring finfish community in the GOM also displayed increases in trophic level.

\begin{figure}

{\centering \includegraphics[width=1\linewidth]{../../../AppData/Local/Temp/RtmpIjNXEs/file2644575b2054} 

}

\caption{Fish community functional traits (growth rate) in New England based on Fall (red) and Spring (blue) survey data with an increasing long-term trend (orange). Dashed lines represent the long-term annual mean for each survey.}\label{fig:traits-k-ne}
\end{figure}

\begin{figure}

{\centering \includegraphics[width=1\linewidth]{../../../AppData/Local/Temp/RtmpIjNXEs/file26441fa2ba1} 

}

\caption{Fish community functional traits (trophic level) in New England based on Fall (red) and Spring (blue) survey data with increasing long-term trend (orange) and increasing recent trend (dark orange). Dashed lines represent the long-term annual mean for each survey.}\label{fig:traits-tl-ne}
\end{figure}

\newpage

Increasing \href{https://noaa-edab.github.io/catalog/exp_n.html}{Adult fish diversity} (Fig. \ref{fig:exp-n-ne}) and changes in functional traits such as the mean trophic level suggest that fish communities have changed from a historic baseline. Long-term trends in biomass and functional traits would not be expected in a stable system. However, because the biomass of functional groups of fish has remained relatively constant over time, the system appears able to adapt to change.

\subsubsection{Implications}\label{implications-5}

The long-term changes in GB suggest an increasingly productive ecosystem from zooplankton to higher trophic levels, but this is occurring simultaneously with low productivity of managed species. Over the same time period, fisheries on GB have experienced cyclic changes in revenue and increased reliance on a single species, sea scallops. Coupled with demographic changes in fisher populations and a decline in the number of New England fleets, this indicates that fisheries utilizing GB may have a lower capacity to adapt to the changing ecosystem. This lower adaptive capacity with a significant shift from baseline conditions in the ecosystem indicate that both the fishery and ecosystem are currently not stable.

Within the GOM, managed species continue to have low productivity, while there are long-term increases in large zooplankton and planktivores. Cyclic changes in zooplankton communities may make the impact of these changes unpredictable. As these changes in the ecosystem occur, an increasing proportion of total revenue is generated by the lobster fishery. This increased reliance on a single species reduces the regions' ability to adapt to changes in resource availability and the environment. For these reasons both the GOM ecosystem and its fisheries are considered not stable.

\newpage

\section{Risks to Meeting Fishery Management Objectives}\label{risks-to-meeting-fishery-management-objectives}

\subsection{Climate and Ecosystem Change}\label{climate-and-ecosystem-change}

\subsubsection{Risks to managing spatially}\label{risks-to-managing-spatially}

Shifting species distributions, or (changes in spatial extent or center of distribution), alter both species and fishery interactions. In particular, shifting species distributions can affect expected management outcomes when spatial allocations and bycatch measures are based on historical fish and protected species distributions. Species availability to surveys can also change as distributions shift within or outside of survey footprints, complicating the interpretation of survey trends.

Coastwide indicators are reviewed in this section to evaluate spatial change throughout the Northeast US shelf. Indicators are identical between the Mid-Atlantic and New England reports.

\paragraph{Indicators: Fish and protected species distribution shifts}\label{indicators-fish-and-protected-species-distribution-shifts}

As noted in the \hyperref[implications]{Seafood Production Implications section}, the combined center of \href{https://noaa-edab.github.io/catalog/species_dist.html}{distribution} for 48 Northeast Shelf commercially or ecologically important fish species continues to show movement towards the northeast and generally into deeper water (Fig. \ref{fig:species-dist}). An analysis of recreational landings data from 2002 to 2019 found evidence of distribution shifts for several \href{https://noaa-edab.github.io/catalog/species_dist.html}{highly migratory species}, including sharks, billfish and tunas.

\href{https://noaa-edab.github.io/catalog/habitat_diversity.html}{Habitat model-based species richness} suggests shifts of both cooler and warmer water species to the northeast. Similar patterns have been found for \href{https://noaa-edab.github.io/catalog/cetacean_dist.html}{marine mammals}, with multiple species shifting northeast between 2010 and 2017 in most seasons (Fig. \ref{fig:protectedspp-dist-shifts}).

Megabenthos center of gravity shows a short-term northward and eastward trend in spring (Fig. \ref{fig:megabenthosshifts}). Megabenthos are large, non-federally-managed benthic invertebrates sampled by scallop dredge, otter trawl, and the Campbell grab. These include crabs, decapods, and sea stars, which are often prey for many managed species.

\begin{figure}

{\centering \includegraphics[width=1\linewidth]{../../../AppData/Local/Temp/RtmpIjNXEs/file264447d9237c} 

}

\caption{Aggregate species distribution metrics for species in the Northeast Large Marine Ecosystem: along shelf distance with increasing trend (orange), and depth with decreasing trend indicating deeper water (purple).}\label{fig:species-dist}
\end{figure}

\begin{figure}

{\centering \includegraphics[width=1\linewidth]{../../../AppData/Local/Temp/RtmpIjNXEs/file26443c781c30} 

}

\caption{Direction and magnitude of core habitat shifts, represented by the length of the line of the seasonal weighted centroid for species with more than 70 km difference between 2010 and 2017 (tip of arrow).}\label{fig:protectedspp-dist-shifts}
\end{figure}

\begin{figure}

{\centering \includegraphics[width=1\linewidth]{../../../AppData/Local/Temp/RtmpIjNXEs/file264471cb23db} 

}

\caption{Eastward (top) and northward (bottom) shifts in the center of gravity for megabenthos species on the Northeast U.S. Shelf in fall (left) and spring (right), with recent increasing trend (orange) for spring eastward center of gravity.}\label{fig:megabenthosshifts}
\end{figure}

\paragraph{Drivers:}\label{drivers}

Mobile populations shift distributions to maintain suitable habitat and prey fields, possibly expanding ranges if new suitable habitat exists. Changes in managed species distribution is partially related to the \href{https://noaa-edab.github.io/catalog/forage_index.html}{distribution of forage biomass}. Since 1982, the fall center of gravity of forage fish (20 species combined) has moved to the north and east (Fig. \ref{fig:forageshifts}). Spring forage fish center of gravity has moved northward but without an eastward trend. Some of the whale and dolphin distribution shifts (Fig. \ref{fig:protectedspp-dist-shifts}) are likely in response to these forage fish shifts. \href{https://noaa-edab.github.io/catalog/zooplankton_index.html\#key-results-and-visualizations-4}{Small copepods}, widespread prey of many larval and juvenile fish, show a similar shift in center of gravity as forage fish, to the north and east in the fall, as well as northward in spring.

\begin{figure}

{\centering \includegraphics[width=1\linewidth]{../../../AppData/Local/Temp/RtmpIjNXEs/file264479db5ed4} 

}

\caption{Eastward (top) and northward (bottom) shifts in the center of gravity for *Calanus finmarchicus* on the Northeast U.S. Shelf in fall (left) and spring (right)}\label{fig:calfin-cog}
\end{figure}
\begin{figure}

{\centering \includegraphics[width=1\linewidth]{../../../AppData/Local/Temp/RtmpIjNXEs/file26446e481379} 

}

\caption{Eastward (top) and northward (bottom) shifts in the center of gravity for 20 forage fish species on the Northeast U.S. Shelf in fall (left) and spring (right), with increasing trends (orange) for fall eastward and northward and spring northward center of gravity.}\label{fig:forageshifts}
\end{figure}

\newpage

In contrast, \href{https://noaa-edab.github.io/catalog/benthos_index.html}{macrobenthos} center of gravity has shifted west and south in the spring (Fig. \ref{fig:macrobenthosshifts}). Macrobenthos are small bottom-dwelling invertebrates including polychaete worms, small crustaceans, bivalves (non-commercial), gastropods, nemerteans, tunicates, cnidarians, brittle stars, sea cucumbers, and sand dollars, and are prey for many managed species. \href{https://noaa-edab.github.io/catalog/zooplankton_index.html\#key-results-and-visualizations-4}{Large copepods} (including \emph{Calanus finmarchicus}) and euphausiids do not have long-term trends in their centers of gravity (Fig. \ref{fig:lgcopeall-cog}) (Fig. \ref{fig:euph-cog}), but small copepods show shifts eastward and northward. Some targeted species distributions may shift in response to these shifts in forage, copepod, and macrobenthos distributions.

\begin{figure}

{\centering \includegraphics[width=1\linewidth]{../../../AppData/Local/Temp/RtmpIjNXEs/file2644283a590c} 

}

\caption{Eastward (top) and northward (bottom) shifts in the center of gravity for macrobenthos species on the Northeast U.S. Shelf in fall (left) and spring (right), with decreasing trends (purple) for spring eastward and northward center of gravity.}\label{fig:macrobenthosshifts}
\end{figure}

\begin{figure}

{\centering \includegraphics[width=1\linewidth]{../../../AppData/Local/Temp/RtmpIjNXEs/file26443f5b5b2a} 

}

\caption{Eastward (top) and northward (bottom) shifts in the center of gravity for large copepod species on the Northeast U.S. Shelf in fall (left) and spring (right) with a  recent decreasing trend (purple) for spring eastward center of gravity.}\label{fig:lgcopeall-cog}
\end{figure}

\begin{figure}

{\centering \includegraphics[width=1\linewidth]{../../../AppData/Local/Temp/RtmpIjNXEs/file264454593090} 

}

\caption{Eastward (top) and northward (bottom) shifts in the center of gravity for Euphausiid species on the Northeast U.S. Shelf in fall (left) and spring (right). }\label{fig:euph-cog}
\end{figure}

Ocean temperatures influence the distribution, seasonal timing, and productivity of managed species (see sections below). The Northeast US shelf, including the Mid-Atlantic, has experienced a continued warming trend for both the \href{https://noaa-edab.github.io/catalog/long_term_sst.html}{long term annual sea surface} (Fig. \ref{fig:long-term-sst}) and \href{https://noaa-edab.github.io/catalog/seasonal_oisst_anom.html}{seasonal surface} and \href{https://noaa-edab.github.io/catalog/bottom_temp_model_anom.html}{bottom temperature}. However, 2025 surface and bottom temperatures were near normal to cooler than normal conditions in all seasons in the MAB (see also the \hyperref[highlights]{2025 Highlights section}).

\begin{figure}

{\centering \includegraphics[width=1\linewidth]{../../../AppData/Local/Temp/RtmpIjNXEs/file264437b77893} 

}

\caption{Seasonal Optimum Interpolation Sea Surface Temperature (OISST) anomaly by season for the MAB, with increasing trends (orange).}\label{fig:seasonal-oisst-anom}
\end{figure}

\begin{figure}

{\centering \includegraphics[width=1\linewidth]{../../../AppData/Local/Temp/RtmpIjNXEs/file26442510188e} 

}

\caption{GLORYS (black) and debiased ROMS (red) seasonal bottom temperature anomaly in the MAB, with increasing trends (orange).}\label{fig:bottom-temp-anom}
\end{figure}

Species suitable habitat can expand or contract when changes in temperature and major oceanographic conditions alter distinct water mass habitats. The variability of the Gulf Stream is a major driver of the predominant oceanographic conditions of the Northeast U.S. continental shelf. As the \href{https://noaa-edab.github.io/catalog/gsi.html}{Gulf Stream} had become less stable and was shifting northward until a recent shift in 2023. Since then, the Gulf Stream has been closer to the long-term average, and the supply of Labrador Slope Water to the Northwest Atlantic Shelf has increased. These changes are linked to some of the cooler water temperatures observed in 2024 and 2025 and the composition of the source water entering the Gulf of Maine through the Northeast Channel (see \hyperref[highlights]{2025 Highlights}).

\begin{figure}

{\centering \includegraphics[width=1\linewidth]{../../../AppData/Local/Temp/RtmpIjNXEs/file26446e39607} 

}

\caption{Northeast U.S. annual sea surface temperature (SST, black), with increasing trend (orange).}\label{fig:long-term-sst}
\end{figure}

\begin{figure}

{\centering \includegraphics[width=1\linewidth]{../../../AppData/Local/Temp/RtmpIjNXEs/file264441737381} 

}

\caption{Index representing changes in the location of the western (between 64 and 55 degrees W) Gulf Stream north wall (black). Positive values represent a more northerly Gulf Stream position, with increasing trend (orange).}\label{fig:GSI}
\end{figure}

Changes in ocean temperature and circulation alter habitat features such as the Mid-Atlantic Bight \href{https://noaa-edab.github.io/catalog/cold_pool.html}{Cold Pool}, a band of relatively cold near-bottom water present from spring to fall over the northern MAB. The cold pool represents essential fish spawning and nursery habitat, and affects fish distribution and behavior. The cold pool has been getting warmer and its areal extent has been shrinking over time (Fig. \ref{fig:cold-pool-size}). In 2025, however, the cold pool temperature index and extent were above the long-term average, likely due to the influx of Labrador Slope and Scotian Shelf waters into the system. Mobile target species that track a preferred temperature range can show increased interannual variability in their distributions as regional temperatures fluctuate from record warms to average over short periods of time.

\begin{figure}

{\centering \includegraphics[width=1\linewidth]{../../../AppData/Local/Temp/RtmpIjNXEs/file26447d406895} 

}

\caption{Seasonal cold pool mean temperature (top) and spatial extent index (bottom), based on bias-corrected ROMS-NWA (open circles) and GLORYS (closed circles), with declining trends (purple).}\label{fig:cold-pool-size}
\end{figure}

\paragraph{Future Considerations}\label{future-considerations}

Distribution shifts caused by changes in thermal habitat and ocean circulation are likely to continue as long as long-term trends persist. Episodic and short-term events (see \href{https://noaa-edab.github.io/catalog/observation_synthesis_2024.html}{2024 Highlights} and \hyperref[highlights]{2025 Highlights}) may increase variability in the trends, however species distributions are unlikely to reverse to historical ranges in the short term. Increased mechanistic understanding of distribution drivers is needed to better understand future distribution shifts: species with high mobility or short lifespans react differently from immobile or long-lived species.

\href{https://noaa-edab.github.io/catalog/MOM6_forecasts.html}{MOM6 decadal oceanographic forecasts} suggest a tendency towards near-normal temperatures over the next decade due to decadal variability in regional circulation. 2026 seasonal forecasts show a high probability of below average surface and bottom temperatures in the winter months. Forecast uncertainty is higher during the spring and summer seasons, and above average conditions are predicted for the fall. These forecasts will continue to be evaluated to determine how well they are able to predict episodic and anomalous events that are outside of the long-term patterns.

Adapting management to changing stock distributions and dynamic ocean processes will require continued monitoring of populations in space and time while evaluating management measures against a range of possible future spatial distributions. The upcoming Climate Vulnerability Assessment 2.0 will also be incorporating MOM6 output and forecasts to help predict changes in species distributions and quantify species exposure to predicted future change. Processes like the \href{https://www.mafmc.org/e3cg}{East Coast Coordination Group} and the HMS Climate Vulnerability Assessment can help coordinate management.

\newpage

\subsubsection{Risks to managing seasonally}\label{risks-to-managing-seasonally}

The effectiveness of seasonal management actions (fishing seasons or area opening/closing periods) depends on a proper alignment with the seasonal life cycle events (phenology) of fish stocks (e.g., migration and spawning timing). If not accounted for, changes in the timing of these biological cycles can reduce the effectiveness of seasonal management measures. The timing of seasonal patterns can also change the interactions between fisheries and non-target species thus influencing bycatch and the availability of species to surveys.

\paragraph{Indicators: Timing shifts}\label{indicators-timing-shifts}

Indicators of phenological changes in fish populations require regular sampling and observations, and therefore a limited number of these indicators are currently available. One indicator shows shifts in \href{https://noaa-edab.github.io/catalog/spawn_timing.html}{spawning timing} of haddock and yellowtail flounder. Spawning of both haddock stocks occurred earlier in the year, as indicated by more resting (post-spawning) stage fish in recent years compared to earlier in the time series (Fig. \ref{fig:spawntiming}). The high percentage of northern stock (Cape Cod/GOM) yellowtail flounder females in the resting maturity stage shown earlier in the time series is reflective of spring surveys sampling them well before spawning, which peaks in June for the northern stock. More recently, the females are much closer to spawning, indicating that yellowtail flounder are spawning earlier in the year. Similarly, increased catch of post-spawning fish in Southern New England, indicates that the peak spawning of the southern stock has also shifted to earlier in the year. Yellowtail flounder spawning is related to bottom temperature, week of year, and decade sampled for each of the three stocks. Changes to spawning times could impact the survival of early life stages of fish, subsequently affecting the larger population size, health, and market value.

\begin{figure}

{\centering \includegraphics[width=1\linewidth]{../../../AppData/Local/Temp/RtmpIjNXEs/file2644497f7e19} 

}

\caption{Percent resting stage (non-spawning) mature female fish (black) from spring NEFSC bottom trawl survey with significant increases (orange) and decreases (purple) from two haddock and three yellowtail flounder stocks: CC = Cape Cod Gulf of Maine, GOM = Gulf of Maine, GB = Georges Bank, SNE = Southern New England.}\label{fig:spawntiming}
\end{figure}

\newpage

\href{https://noaa-edab.github.io/catalog/timing_shifts.html}{Migration timing} of some tuna and large whale migrations has changed. An analysis of recreational fishing data between 2019 and 2022 identified multiple shifts in important HMS species. For example, Bigeye tuna were caught 50 days earlier; small and large bluefin tuna were caught 38 and 80 days earlier respectively in Massachusetts; and blue marlin in New York were caught 27 days earlier. A separate analysis of acoustic telemetry data predicted delayed departure of southward-migrating sharks from the northeast region under future sea surface temperatures. These results are further supported by the Atlantic Highly Migratory Species Climate Vulnerability Assessment, which found that 57 of 58 highly migratory species and stocks have high or very high potential to shift distributions. In Cape Cod Bay, peak spring habitat use by right and humpback whales has shifted 18-19 days later over time.

Understanding whether seasonal patterns are changing for stocks requires regular observations throughout the year. For example, baseline work on \href{https://noaa-edab.github.io/catalog/cetacean_acoustic.html}{cetacean presence in Southern New England} shows different seasonal use patterns for whale and dolphin species. Despite the importance of understanding seasonal patterns, we have few indicators that directly assess timing shifts of species. We plan on incorporating more indicators of timing shifts and phenology in future reports.

\paragraph{Drivers:}\label{drivers-1}

The drivers of timing shifts in managed stocks are generally coupled to shifts in environmental or biological conditions, since these can result in changes in habitat quality or food availability within the year. Changes in the timing of fall phytoplankton blooms and seasonal shifts in zooplankton communities are indicators of changes in seasonal food availability to stocks.

Along with the overall warming trends in New England, ocean summer conditions have been lasting longer (Fig. \ref{fig:transition}) due to the earlier \href{https://noaa-edab.github.io/catalog/trans_dates.html}{transition} from cool spring conditions to warm summer conditions and the later transition from warm summer conditions to cooler fall temperatures. These transition dates relate how daily temperatures compare to the seasonal norm. Changes in the timing of seasonal environmental cycles can alter biological processes (migrations, spawning, etc.) that are triggered by seasonal events.

\begin{figure}

{\centering \includegraphics[width=1\linewidth]{../../../AppData/Local/Temp/RtmpIjNXEs/file264460c95d80} 

}

\caption{Ocean summer length in New England (Georges Bank, top; Gulf of Maine, bottom): the annual total number of days between the spring thermal transition date and the fall thermal transition date (black), with an increasing trend (orange). Transition dates are based on sea surface temperatures.}\label{fig:transition}
\end{figure}

\newpage

The Mid-Atlantic Bight \href{https://noaa-edab.github.io/catalog/cold_pool.html}{Cold Pool} is a summer to early fall feature that creates seasonally suitable habitat for many species, including some managed by the NEFMC. Cold pool persistence has decreased, indicating that the duration of the cold pool habitat is shorter compared to the 1960s (Fig. \ref{fig:cold-pool-time-ne}). However, all cold pool indices were near or above the long-term average in 2025 and likely related to the influx of northern waters into the system (see \href{https://noaa-edab.github.io/catalog/observation_synthesis_2024.html}{2024 Highlights} ). A change in the timing of the autumn breakdown of the Cold Pool may impact the recruitment of species that rely on it for seasonal cues and habitat. Southern New England-Mid Atlantic yellowtail flounder recruitment and settlement are related to the strength of the Cold Pool (a factor of extent and persistence). The correlation of pre-recruit settlers to the Cold Pool is thought to represent a bottleneck in yellowtail flounder life history, whereby a local and temporary increase in bottom temperature can negatively impact the survival of settlers. Including the effect of Cold Pool variations on yellowtail recruitment reduced retrospective patterns and improved predictive skill in a stock assessment model. This connection is especially important given the long-term decline in the duration of the Cold Pool.

\begin{figure}

{\centering \includegraphics[width=1\linewidth]{../../../AppData/Local/Temp/RtmpIjNXEs/file26444c5654d7} 

}

\caption{The Mid-Atlantic Bight Cold Pool persistence index based on bias-corrected ROMS-NWA (open circles) and GLORYS (closed circles), with significant long-term decline (purple).}\label{fig:cold-pool-time-ne}
\end{figure}

The seasonal timing of \href{https://noaa-edab.github.io/catalog/chl_pp.html}{phytoplankton} blooms shows a tendency towards an increased fall bloom over time in the GOM and GB, with chlorophyll concentrations significantly increasing in October and November (GB) and January and October (GOM) (Fig. \ref{fig:chl-month-ne}). This increased production at the base of the food web may increase prey availability, and fall blooms in particular have been associated with increased recruitment for species such as haddock.

\begin{figure}

{\centering \includegraphics[width=1\linewidth]{../../../AppData/Local/Temp/RtmpIjNXEs/file26447650cd2} 

}

\caption{Monthly median chlorophyll a concentration time series for Georges Bank (top) and Gulf of Maine (bottom), with increasing trends (orange). }\label{fig:chl-month-ne}
\end{figure}

\paragraph{Future Considerations}\label{future-considerations-1}

Species are reliant on environmental processes to dictate the timing of their behavior (e.g., phytoplankton bloom timing, thermal transition, or the duration of the cold pool). Some changes are episodic and have interannual variability, while others may be shifting away from a historic baseline on the scales of years to decades. Other species may rely on the general seasonal succession of environmental drivers (e.g., the timing of the fall turnover) to cue biological processes, and long-term trends in seasonal transitions are unlikely to reverse in coming years. Thus, timing shifts in migration or spawning may continue. Management actions that rely on effective alignment of fisheries availability and biological processes should continue to evaluate whether prior assumptions on seasonal timings still hold, and new indicators should be developed to monitor timing shifts for stocks.

\newpage

\subsubsection{Risks to setting catch limits}\label{risks-to-setting-catch-limits}

The efficacy of short-term stock projections and rebuilding plans rely on accurate understanding of processes affecting stock growth, reproduction, and natural mortality. These biological processes are often driven by underlying environmental change. If ignored, environmental change may increase the risk that established stock-level biological reference points no longer reflect the current population and increase projection uncertainty, both of which can contribute to quota misspecification.

\paragraph{Indicators: Fish productivity and condition shifts}\label{indicators-fish-productivity-and-condition-shifts}

Indicators of \href{https://noaa-edab.github.io/catalog/productivity_anomaly.html}{fish productivity} derived from observations (surveys) or models (stock assessments) show variability over the time series. Since 2020, fish productivity has been below the long-term average productivity (derived from NEFSC bottom trawl survey), and 2025 was below average for all managed species (Fig. \ref{fig:productivity-anomaly-ne}). A similar analysis based on stock assessment model outputs shows a decline in productivity over the time series with relatively high productivity in the 1990s and relatively low productivity in the 2000's (recruitment per spawning stock biomass anomaly). Fish productivity can be affected by parental condition, environmental conditions, timing and availability of prey for recruits, as well as retention of recruits within favorable habitat. In years where offshore advection is high in a depth range and month when a fish species spawns, fish productivity and recruitment may be reduced. Other signs of changing productivity in New England are the declines in \href{https://noaa-edab.github.io/catalog/seabird_ne.html}{common tern chicks} per nest (Fig. \ref{fig:seabird-ne-productivity}) and continued low overall \href{https://noaa-edab.github.io/catalog/gom_salmon.html}{Atlantic salmon} abundance (Fig. \ref{fig:salmon}) despite short-term increases in return rates and salmon numbers.

\begin{figure}

{\centering \includegraphics[width=1\linewidth]{../../../AppData/Local/Temp/RtmpIjNXEs/file264460ab6e1b} 

}

\caption{Fish productivity measures. Top: Small-fish-per-large-fish survey biomass anomaly for species managed by the New England Fishery Management Council. Bottom: assessment recruitment per spawning stock biomass anomaly for stocks managed by the New England Fishery Management Council. The summed anomaly across species is shown by the black line, drawn across all years with the same number of stocks analyzed.}\label{fig:productivity-anomaly-ne}
\end{figure}

\begin{figure}

{\centering \includegraphics[width=1\linewidth]{../../../AppData/Local/Temp/RtmpIjNXEs/file264474a129e4} 

}

\caption{Common tern productivity - the number of fledged chicks per nest - at seven Gulf of Maine colonies managed by the National Audubon Society's Seabird Restoration Program, with a significant long-term decline (purple).}\label{fig:seabird-ne-productivity}
\end{figure}

\begin{figure}

{\centering \includegraphics[width=1\linewidth]{../../../AppData/Local/Temp/RtmpIjNXEs/file264468371887} 

}

\caption{Percent return rate (top) and abundance (bottom) of Atlantic salmon returns to Gulf of Maine rivers since 1972 and return rates for two sea winter returns from hatchery smolt stocking in the Penobscot River. Long-term decreasing trend for percent return rate in purple.}\label{fig:salmon}
\end{figure}

The health of individual fish (i.e., fish condition) can contribute to population productivity through improved growth, reproduction, and survival. \href{https://noaa-edab.github.io/catalog/condition.html}{Fish condition} in the Gulf of Maine and Georges Bank regions were generally high to very high prior to 2000, low to very low from 2001-2010 (concurrent with declines in fish productivity, Fig. \ref{fig:productivity-anomaly-ne}), and mixed since 2011. In 2025, fish condition was below or close to average for most species on both Georges Bank and in the Gulf of Maine (Fig. \ref{fig:ne-cf}). Preliminary analyses show that years dominated by small copepods and warmer spring temperatures may improve fish condition for Atlantic mackerel and butterfish. Similar environmental drivers may be important to other species.

\begin{figure}

{\centering \includegraphics[width=1\linewidth]{../../../AppData/Local/Temp/RtmpIjNXEs/file2644730115c9} 

}

\caption{Condition factor for fish species in New England (top: Georges Bank, GB; bottom: Gulf of Maine, GOM) based on fall NEFSC bottom trawl survey data. No survey was conducted in 2020.}\label{fig:ne-cf}
\end{figure}

\paragraph{Drivers}\label{drivers-2}

Fish productivity and condition are the cumulative effects of physiological, ecological, and environmental factors. Major factors include increased metabolic demands from increasing temperature and changes in the availability and quality of prey. Long-term environmental trends and episodic extreme temperatures, ocean acidification, and low oxygen events represent multiple stressors that can affect growth rates, reproductive success, recruitment, and cause mortality.

\subparagraph{Biological Drivers: Forage quality and abundance}\label{biological-drivers-forage-quality-and-abundance}

Management should account for energetic links between prey and predators, as shifts in forage quality and abundance directly alter the health, productivity, and movement of managed and protected species. The total energy available to higher trophic level predators is determined by the mass and energy density (ED) of prey. Protecting this forage base is essential for maintaining overall ecosystem function and continued stock productivity and condition.

Forage \href{https://noaa-edab.github.io/catalog/energy_density.html}{energy content} fluctuates based on growth, reproduction, and environmental productivity. High-energy New England species include alewife, Atlantic mackerel, and Atlantic herring. Alewife provide the highest ED in the GOM during the fall. Atlantic mackerel show higher abundance and ED in the GOM than on GB during fall. Atlantic herring offer a consistent year-round energy source, though values vary between spring and fall spawning groups. Butterfish abundance has increased over the last five years in both regions during the fall, providing an additional high-energy prey option.

Moderate-energy species, including longfin squid, shortfin squid, and silver hake, provide a stable but lower ED food supply. Squid abundance is generally lower in the GOM than on GB. Silver hake remain highly abundant in the GOM with stable ED across spring and fall. Northern sand lance offer intermediate energy but are only available in the spring before burying in the seafloor to overwinter.

Declining prey energy density creates significant risks for both forage and predator stocks. In prey species like silver hake, lower energy reserves can reduce spawning success and recruitment. For predators, including managed species such as goosefish and spiny dogfish, lower-quality prey leads to poorer physical health and reduced reproductive output.

Shifts in forage abundance directly influence managed species productivity. While New England fall \href{https://noaa-edab.github.io/catalog/forage_index.html}{forage biomass} remains stable, long-term increases are observed in the spring GOM. Biomass peaked during the 1980s in the fall. Increased spring GOM (Fig. \textbackslash ref\{fig:energy-density-ne\}) forage biomass may improve fish health and reproductive output during spawning seasons when energy reserves are typically low. However, this benefit may be offset by lower prey energy densities, particularly during periods of higher water temperatures when predator metabolic demands increase.

\begin{figure}

{\centering \includegraphics[width=1\linewidth]{../../../AppData/Local/Temp/RtmpIjNXEs/file26442ba6681} 

}

\caption{Energy density (mean and standard deviation) of eight forage species from NEFSC bottom trawl surveys by season and year for the GOM. Symbol size represents abundance (mean kg/tow) estimated from bottom trawl survey tows in GB (top) and the GOM (bottom). }\label{fig:energy-density-ne}
\end{figure}

\begin{figure}

{\centering \includegraphics[width=1\linewidth]{../../../AppData/Local/Temp/RtmpIjNXEs/file264478ea7901} 

}

\caption{Forage fish index in GB (top) and GOM  (bottom, significant long-term increase for spring in orange) for spring (blue) and fall (red) surveys. Index values are relative to the maximum observation within a region across surveys.}\label{fig:foragebio-ne}
\end{figure}

\href{https://noaa-edab.github.io/catalog/benthos_index.html}{Benthic invertebrates} are extremely important forage for some managed species (e.g., flatfish, juvenile cod and haddock) Macrobenthos indices show long term declines in spring. In contrast, megabenthos indices show long-term increases during the fall in both GB and GOM (Fig. \ref{fig:benthos-ne}). Fish productivity may be negatively impacted in recent years for fish such as flounders and juvenile fish that target macrobenthos such as small crustaceans and polychaetes in the spring, and positively impacted for fish such as larger skates, hakes and gadids that target megabenthos such as crabs.

\subparagraph{Biological Drivers: Lower trophic levels}\label{biological-drivers-lower-trophic-levels}

Phytoplankton are the foundation of the marine food web and are the primary food source for zooplankton and filter feeders such as shellfish. Multiple environmental and oceanographic drivers affect the abundance, composition, spatial distribution, and productivity of phytoplankton. While changes in phytoplankton productivity could affect fish productivity (including forage), there is no clear long-term trend in New England total primary production (Fig. \ref{fig:totpp-ne}).

New England \href{https://noaa-edab.github.io/catalog/zoo_abundance_anom.html}{zooplankton abundance} is shifting in ways that could impact fish condition and marine mammal prey availability. In the Gulf of Maine (GOM), increased small-bodied copepods and euphausiids are linked to improved condition in species like Atlantic mackerel, and baleen whales (humpback, sei, and fin) may benefit from long-term increases in prey availability, although euphausiid biomass has been recently high variable (Fig. 83). Conversely, large-bodied copepods in Georges Bank (GB) have declined recently. Zooplankton energy density varies by season and location, with high-energy large copepods peaking from April through June. Since 2023, \href{https://noaa-edab.github.io/catalog/zoo_community.html}{zooplankton communities} have reverted to compositions similar to pre-1990 and 2000-2011 periods; research is currently underway to determine the drivers and management implications of these shifts.

\begin{figure}

{\centering \includegraphics[width=1\linewidth]{../../../AppData/Local/Temp/RtmpIjNXEs/file26441f4163da} 

}

\caption{Changes in spring (blue) and fall (red) benthos abundance in New England for megabenthos (top) and macrobenthos (bottom), with significant long-term increasing (orange), short-term increasing (dark orange), and long-term decreasing (purple) trends. }\label{fig:benthos-ne}
\end{figure}

\emph{Calanus finmarchicus} abundance has declined in the GOM following a 2008 shift in oceanographic conditions, which poses a risk to the critically endangered North Atlantic right whale and key energy link in subarctic ecosystems. This lipid-rich copepod can comprise 71\% of the total zooplankton biomass. Observations in the Wilkinson Basin indicate that the spring and summer abundance and biomass of Calanus in 2024 was comparable to 2005. However, \href{https://noaa-edab.github.io/catalog/wbts_zoo.html}{late-stage abundance} has declined 64\% in fall and 71\% in winter. Consequently, overall mesozooplankton biomass in 2024 was only 27\% of 2005 levels (Fig. \ref{fig:wbts-zoo}).

The seasonal differences in \emph{Calanus} are driven by five factors:
1. Late winter and early spring phytoplankton levels control reproductive output.
2. Source water origin determines supply, with higher concentrations in \href{https://noaa-edab.github.io/catalog/wbts_zoo.html}{Scotian/Labrador shelf water} than in warm slope water.
3. Variable currents dictate how \emph{Calanus} is transported and retained in deep basins.
4. Invertebrate predator populations fluctuate based on spring \emph{Calanus} abundance.
5. Higher summer and fall temperatures increase predator metabolic demands and predation pressure.

\begin{figure}

{\centering \includegraphics[width=1\linewidth]{../../../AppData/Local/Temp/RtmpIjNXEs/file264473e84814} 

}

\caption{Changes in three dominant zooplankton (\textit{Calanus finmarchicus}, \textit{Calanus typicus}, and \textit{Pseudocalanus spp}) abundance anomalies in New England for large (top), small (middle) copepods, and Euphausiids (bottom), with significant decreases (short-term, dark purple) in large copepods and long-term increases (orange) in GOM small copepods and GB and GOM euphausiids.}\label{fig:zoopanom-ne}
\end{figure}

\begin{figure}

{\centering \includegraphics[width=1\linewidth]{../../../AppData/Local/Temp/RtmpIjNXEs/file26447fbc3827} 

}

\caption{Mesozooplankton biomass phenology at Wilkinson Basin Time Series station: 2005 to 2025. Total dry weight from 200m mesh vertical ring net tows. Fitted line shows GAM predictions with 95\% confidence interval (shaded). Circles:  2004 to 2010; triangles: 2011 to 2022; 2023, 2024, and 2025 shown as separate symbols (see legend). Vertical lines denote season boundaries. Data for 2025 incomplete. Climatology GAM prediction calculated with a year of 2012; prediction is significantly different from random pattern. triangles: 2011 to 2024; squares: 2025.}\label{fig:wbts-zoo}
\end{figure}

\subparagraph{Environmental Drivers}\label{environmental-drivers}

Fish production can also be directly related to the prevailing environmental conditions by altering metabolic (growth) and reproductive processes. Many species possess thermal tolerances and can experience stressful or lethal conditions if temperatures exceed certain levels. Extreme temperature at both the \href{https://noaa-edab.github.io/catalog/seasonal_oisst_anom.html}{surface} (Fig. \ref{fig:long-term-sst}) and \href{https://noaa-edab.github.io/catalog/bottom_temp_model_anom.html}{bottom} can exceed \href{https://noaa-edab.github.io/catalog/thermal_habitat_gridded.html}{thermal tolerance} limits for some fish. For example, 2012 had among the warmest surface and bottom temperatures (GB) in New England. A large proportion of the Georges Bank and Mid-Atlantic regions had bottom temperatures above the 15℃ thermal tolerance for most groundfish, with some days in the Mid-Atlantic exceeding the 24℃ potential mortality limit (Fig. \ref{fig:therm-hab-persist-2024-ne}).

Cooler ocean temperatures prevented \href{https://noaa-edab.github.io/catalog/heatwave.html}{marine heatwaves} in the Gulf of Maine and Georges Bank during 2025. Instead, Georges Bank experienced three surface and two bottom marine cold spells, which are extreme cooling events below the 90th percentile. The location, duration, and timing of \href{https://noaa-edab.github.io/catalog/cold_spell.html}{cold spells} can affect the productivity of temperature-sensitive species. The most significant surface event occurred in November, ranking as the 11th strongest on record, while a notable bottom cold spell beginning August 11th reached peak intensity on September 15th and may be ongoing. Another bottom cold spell on the Bank persisted for 71 days starting in early February.

The Gulf of Maine recorded five surface and three bottom marine cold spells in 2025. A major surface event began February 6th and lasted 42 days, with sea surface temperatures averaging 4.50 °C---nearly 1 °C below the 2016-2025 average. Additionally, the seventh strongest surface cold spell on record occurred in April, lasting 37 days. Bottom conditions in the Gulf of Maine were similarly impacted by three cold spells, including the fifth strongest on record. This event began in December 2024 and lasted 83 days, with bottom temperatures averaging 7.5 °C. This represented a cooling of more than 1 °C compared to the 2016-2025 average.

Lower ocean temperatures near long-term averages will affect species differently across the region. While cold-water species like cod may benefit from these conditions, warm-water species such as black sea bass are unlikely to see positive effects. This variability in regional cooling highlights the need for management to account for shifting species distributions and productivity.

\begin{figure}

{\centering \includegraphics[width=1\linewidth]{../../../AppData/Local/Temp/RtmpIjNXEs/file26442b425762} 

}

\caption{Marine cold spells based on detrended GLORYS12V1 bottom temperature from October 2024 though October 2025 in the GOM. Cold spells defined by events where observered temperature (black) exceeds  a threshold (green) based on the climatology (blue). Blue shaded areas show the cold spell events with the largest event in dark blue. }\label{fig:mcs-bottom-gom}
\end{figure}

\begin{figure}

{\centering \includegraphics[width=1\linewidth]{../../../AppData/Local/Temp/RtmpIjNXEs/file26443aa15d7d} 

}

\caption{The number of days in 2025 where bottom temperature exceeds 15 degrees (left) and 24 degrees (right) based on the GLORYS 1/12 degree grid.}\label{fig:therm-hab-persist-2024-ne}
\end{figure}

The newly-developed \href{https://noaa-edab.github.io/catalog/advection.html}{advection index} (Fig. \ref{fig:advection-index-ne}) shows total transport of water onto and off the continental shelf and can be linked to the survival of early life stages of fish and invertebrates. Long-term trends in New England show increased onshelf movement of mid-layer and bottom waters in June, which could increase retention of some species. Further study is needed on the species level to link spawning timing and larval periods to the advection index at the corresponding depth and month.

\begin{figure}

{\centering \includegraphics[width=1\linewidth]{../../../AppData/Local/Temp/RtmpIjNXEs/file26443b7f1885} 

}

\caption{Net advection in June across the Southern New England and Georges Bank continental shelf break within 3 depth bands, with increasing long-term trends in the middle and bottom depths (orange).}\label{fig:advection-index-ne}
\end{figure}

\href{https://noaa-edab.github.io/catalog/ocean_acidification.html}{Ocean acidification} risks vary among species and include reduced survival, growth, reproduction, and productivity,reached record levels in 2024, though were moderated in 2025. Atlantic sea scallop and longfin squid faced high OA risk in Long Island Sound and the New Jersey shelf during the summers of 2016, 2018, 2019, 2023, and 2024, with 2024 marking the highest risk recorded since 2007. By 2025, risk levels decreased but still exceeded biological sensitivity limits for scallops on the New Jersey outer shelf in \href{https://noaa-edab.github.io/catalog/ocean_acidification.html}{spring} and reached sensitivity limits for longfin squid in nearshore New Jersey waters during summer. These risks are heightened by cold-water CO2 absorption and the movement of high-CO2 water masses. While 2025 bottom temperatures remained as cool as 2024, higher salinity indicated a shift in water mass composition that resulted in lower overall OA risk compared to the previous two years.

\begin{figure}

{\centering \includegraphics[width=1\linewidth]{../../../AppData/Local/Temp/RtmpIjNXEs/file26445ffa5b6e} 

}

\caption{Locations where bottom aragonite saturation state ($\Omega_{Arag}$; summer only: June-August) were at or below the laboratory-derived sensitivity level for Atlantic sea scallop (left panel) and longfin squid (right panel) for the time periods 2007-2022 (dark cyan), 2023 only (magenta) and 2024 only (cyan). Gray circles indicate locations where bottom $\Omega_{Arag}$ values were above the species specific sensitivity values..}\label{fig:mab-oa}
\end{figure}

Biological and oceanographic processes can affect the amount of oxygen present in the water column. During low oxygen (hypoxic) events, species' growth is negatively affected and very low oxygen can result in mortality. In 2025, \href{https://ondeckdata.com/database/soe_hypoxic_days_2025.html}{aggregated demersal DO observations} collected by a variety of programs were examined simultaneously. These programs include glider deployments, fishery-independent surveys, and cooperative ocean observing efforts aboard commercial fishing vessels. Coastal hypoxia was observed in Narragansett Bay in September and October where water temperatures were warm and stagnant. There were no reports of mass mortality events from the fishing industry. The duration and extent of hypoxic events is being monitored, but long-term shelf-wide observations are not yet available. However, \href{https://noaa-edab.github.io/catalog/observation_synthesis_2023.html}{hypoxic events} were detected in Cape Cod Bay in 2019 and 2020 and off the coast of New Jersey in 2023 and were potentially responsible for fish, lobster, and crab \href{https://sebsnjaesnews.rutgers.edu/2023/12/rutgers-scientists-observe-unusual-ocean-conditions-possibly-linked-to-mortality-in-marine-life-off-new-jersey/}{mortalities}.

\subparagraph{Drivers: Predation}\label{drivers-predation}

The abundance and distribution of predators can affect both the productivity and mortality rates on managed stocks. Predators can consume managed species or compete for the same resources resulting in increased natural mortality or declining productivity, respectively. The northeast shift in some \href{https://noaa-edab.github.io/catalog/cetacean_dist.html}{highly migratory species} (Fig. \ref{fig:protectedspp-dist-shifts}) indicates a change in the overlap between predators and prey. Since we also observe distribution shifts in both managed and forage species, the effect of changing predator distributions alone is difficult to quantify.

\href{https://noaa-edab.github.io/catalog/grayseal.html}{Gray seals} are fish predators with increasing populations in New England,Recent white shark aggregations have been observed near Cape Cod, however, both gray seals and white sharks are broad generalist feeders that do not generally target commercially-sized managed species. \href{https://noaa-edab.github.io/catalog/hms_stock_status.html}{Stock status} is mixed for Atlantic Highly Migratory Species (HMS) stocks (including sharks, swordfish, billfish, and tunas) occurring throughout the Northeast U.S. shelf. While there are several HMS species considered to be overfished or that have unknown stock status, the population status for some managed Atlantic sharks and tunas is at or above the biomass target, suggesting the potential for robust (or rebuilt) predator populations among these managed species. Stable predator populations suggest stable predation pressure on managed species, but increasing predator populations may reflect increasing predation pressure.

\paragraph{Future Considerations}\label{future-considerations-2}

The processes that control fish productivity and mortality are dynamic, complex, and the result of the interactions between multiple system drivers. There is a real risk that short-term predictions in assessments and rebuilding plans that assume unchanging underlying conditions will not be as effective, given the observed ecological and environmental process changes documented throughout the report. Assumptions for species' growth, reproduction, and natural mortality should continue to be evaluated for individual species. With observations of system-wide productivity shifts of multiple managed stocks, more research is needed to determine whether regime shifts or ecosystem reorganization are occurring, and how this should be incorporated into management

\newpage

\subsection{Other Ocean Uses: Offshore Wind}\label{other-ocean-uses-offshore-wind}

Offshore wind development is active and dynamic throughout the region. The following section reflects the status of projects as of 5 February, 2026.

\subsubsection{Indicators: development timeline, revenue in lease areas, coastal community vulnerability}\label{indicators-development-timeline-revenue-in-lease-areas-coastal-community-vulnerability}

All reported potential offshore wind development status and data are based on BOEM's Offshore Renewable Activities page and projects' Final Environmental Impact Statements. In 2025, the Presidential Memorandum 90 FR 8363 removed existing planning areas and excluded the establishment of additional lease areas.

As of January 2026, 38 offshore \href{https://noaa-edab.github.io/catalog/wind_dev_speed.html}{wind development} leases are under different stages of development in the Northeast (Fig. \ref{fig:wind-dev-cumul}). One project (South Fork Wind Farm) is fully operational and another (Vineyard Wind 1) is partly operational while construction finishes. The southern New England region has two other projects currently under construction (Revolution Wind and Sunrise Wind). Empire Wind and Coastal Virginia Offshore Wind (CVOW) are currently under construction in the New York Bight and Mid-Atlantic Region, respectively, with CVOW expected to start generating power in early 2026.

Construction of these projects during 2025 affected fisheries managed by the {[}Mid-Atlantic{]} {[}New England{]} Fishery Management Council. There are eight additional projects that have Construction and Operations Plan (COP) approvals (three in Southern New England and five in the Mid-Atlantic/New York Bight) that could begin construction in 2026, however, construction schedules are highly uncertain at this time. Seven additional projects have submitted COPs and are pending approval, while the remaining projects are under the site assessment phase and have not submitted COPs to date (Fig. \ref{fig:wind-dev-cumul}).

With the first offshore wind energy projects now under construction and operation, all indicator analyses in this section follow a different reporting format than in previous years. Where previous years reported data for all lease areas, this year we investigate impacts of the six commercial scale projects currently under construction or operation, (i.e., Active Projects: South Fork Wind Farm, Revolution Wind, Sunrise Wind, Empire Wind 1, Vineyard Wind 1, and CVOW-Commercial).

\begin{figure}

{\centering \includegraphics[width=0.9\linewidth]{../../../AppData/Local/Temp/RtmpIjNXEs/file264477b0a4} 

}

\caption{All Northeast Project areas by year construction ends (each project has a 2 year construction period).}\label{fig:wind-dev-cumul}
\end{figure}

\begin{figure}

{\centering \includegraphics[width=1\linewidth]{../../../AppData/Local/Temp/RtmpIjNXEs/file26442b2b7d69} 

}

\caption{Total area proposed for wind development on the Northeast Shelf through 2030.}\label{fig:wind-proposed-dev}
\end{figure}

With the first offshore wind energy projects now under construction and operation, all indicator analyses in this section follow a different reporting format than in previous years. Where previous years reported data for all lease areas, this year we investigate impacts of the six commercial-scale projects currently under construction or operation (i.e., Active Projects: South Fork Wind Farm, Revolution Wind, Sunrise Wind, Empire Wind 1, Vineyard Wind 1, and CVOW-Commercial).

Offshore wind indicators are based on federal logbook data and do not include all data for all fisheries; therefore a complete evaluation of potential offshore wind energy development impacts would need to be supplemented by other data sources. For further information on the utility of the data, see the \href{https://www.fisheries.noaa.gov/resource/data/socioeconomic-impacts-atlantic-offshore-wind-development}{socioeconomic impacts of offshore wind development data reports page}.

Based on federal vessel logbook data, \href{https://noaa-edab.github.io/catalog/wind_revenue.html}{commercial fishery revenue} from trips within Active Projects varied annually from 2008-2024. Maximum annual revenue for the fisheries with the most overlap with Active Projects peaked at over \$8.7 million for the sea scallop fishery, \$1.1 million for monkfish, \$477,000 for skates, \$377,000 for yellowtail flounder, and \$344,000 for Atlantic herring (Fig. 92). Individual groundfish species are more affected on a percentage basis, with up to 13\% of historical annual revenues overlapping with Active Projects for species such as little skate (13\%), barndoor skate (11\%), yellowtail flounder (10\%), and 6\% each for red hake, clearnose skate, and winter skate, respectively (Table \ref{tab:wea-landings-rev-ne}). Future fishery resource overlap with wind leases, especially scallops, may change due to species distribution shifts attributable to climate change and recruitment and larval dispersion pattern changes caused by hydrodynamic flow disruptions from turbine foundations, which could also affect fishery landings/revenue.

\begin{figure}

{\centering \includegraphics[width=1\linewidth]{../../../AppData/Local/Temp/RtmpIjNXEs/file2644e0f1669} 

}

\caption{Revenue of species managed by the New England Fishery Management Council within existing offshore wind lease areas.}\label{fig:wea-spp-rev-ne}
\end{figure}

\global\setlength{\Oldarrayrulewidth}{\arrayrulewidth}

\global\setlength{\Oldtabcolsep}{\tabcolsep}

\setlength{\tabcolsep}{2pt}

\renewcommand*{\arraystretch}{1}



\providecommand{\ascline}[3]{\noalign{\global\arrayrulewidth #1}\arrayrulecolor[HTML]{#2}\cline{#3}}

\begin{longtable}[c]{|p{2.00in}|p{2.00in}|p{2.00in}}

\caption{Top\ 10\ New\ England\ managed\ species\ Landings\ and\ Revenue\ from\ Wind\ Energy\ Areas.\ Skates\ includes\ barndoor,\ winter,\ clearnose,\ smooth,\ little,\ and\ general\ skates\ reported\ in\ logbooks.\ *Less\ than\ a\ maximum\ of\ 50,000\ lb\ was\ reported\ landed\ annually\ in\ wind\ energy\ lease\ areas\ for\ these\ species..}\label{tab:wea-landings-rev-ne}\\

\ascline{1.5pt}{666666}{1-3}

\multicolumn{1}{>{\raggedright}m{\dimexpr 2in+0\tabcolsep}}{\textcolor[HTML]{000000}{\fontsize{9}{9}\selectfont{NEFMC,\ MAFMC,\ and\ ASMFC\ Managed\ Species}}} & \multicolumn{1}{>{\raggedleft}m{\dimexpr 2in+0\tabcolsep}}{\textcolor[HTML]{000000}{\fontsize{9}{9}\selectfont{Maximum\ Percent\ Total\ Annual\ Regional\ Species\ Landings}}} & \multicolumn{1}{>{\raggedleft}m{\dimexpr 2in+0\tabcolsep}}{\textcolor[HTML]{000000}{\fontsize{9}{9}\selectfont{Maximum\ Percent\ Total\ Annual\ Regional\ Species\ Revenue}}} \\

\ascline{1.5pt}{666666}{1-3}\endfirsthead \caption[]{Top\ 10\ New\ England\ managed\ species\ Landings\ and\ Revenue\ from\ Wind\ Energy\ Areas.\ Skates\ includes\ barndoor,\ winter,\ clearnose,\ smooth,\ little,\ and\ general\ skates\ reported\ in\ logbooks.\ *Less\ than\ a\ maximum\ of\ 50,000\ lb\ was\ reported\ landed\ annually\ in\ wind\ energy\ lease\ areas\ for\ these\ species..}\label{tab:wea-landings-rev-ne}\\

\ascline{1.5pt}{666666}{1-3}

\multicolumn{1}{>{\raggedright}m{\dimexpr 2in+0\tabcolsep}}{\textcolor[HTML]{000000}{\fontsize{9}{9}\selectfont{NEFMC,\ MAFMC,\ and\ ASMFC\ Managed\ Species}}} & \multicolumn{1}{>{\raggedleft}m{\dimexpr 2in+0\tabcolsep}}{\textcolor[HTML]{000000}{\fontsize{9}{9}\selectfont{Maximum\ Percent\ Total\ Annual\ Regional\ Species\ Landings}}} & \multicolumn{1}{>{\raggedleft}m{\dimexpr 2in+0\tabcolsep}}{\textcolor[HTML]{000000}{\fontsize{9}{9}\selectfont{Maximum\ Percent\ Total\ Annual\ Regional\ Species\ Revenue}}} \\

\ascline{1.5pt}{666666}{1-3}\endhead



\multicolumn{1}{>{\raggedright}m{\dimexpr 2in+0\tabcolsep}}{\textcolor[HTML]{000000}{\fontsize{9}{9}\selectfont{Little\ Skate}}} & \multicolumn{1}{>{\raggedleft}m{\dimexpr 2in+0\tabcolsep}}{\textcolor[HTML]{000000}{\fontsize{9}{9}\selectfont{7.42}}} & \multicolumn{1}{>{\raggedleft}m{\dimexpr 2in+0\tabcolsep}}{\textcolor[HTML]{000000}{\fontsize{9}{9}\selectfont{12.94}}} \\





\multicolumn{1}{>{\raggedright}m{\dimexpr 2in+0\tabcolsep}}{\textcolor[HTML]{000000}{\fontsize{9}{9}\selectfont{Barndoor\ Skate*}}} & \multicolumn{1}{>{\raggedleft}m{\dimexpr 2in+0\tabcolsep}}{\textcolor[HTML]{000000}{\fontsize{9}{9}\selectfont{11.45}}} & \multicolumn{1}{>{\raggedleft}m{\dimexpr 2in+0\tabcolsep}}{\textcolor[HTML]{000000}{\fontsize{9}{9}\selectfont{11.07}}} \\





\multicolumn{1}{>{\raggedright}m{\dimexpr 2in+0\tabcolsep}}{\textcolor[HTML]{000000}{\fontsize{9}{9}\selectfont{Yellowtail\ Flounder}}} & \multicolumn{1}{>{\raggedleft}m{\dimexpr 2in+0\tabcolsep}}{\textcolor[HTML]{000000}{\fontsize{9}{9}\selectfont{9.01}}} & \multicolumn{1}{>{\raggedleft}m{\dimexpr 2in+0\tabcolsep}}{\textcolor[HTML]{000000}{\fontsize{9}{9}\selectfont{9.61}}} \\





\multicolumn{1}{>{\raggedright}m{\dimexpr 2in+0\tabcolsep}}{\textcolor[HTML]{000000}{\fontsize{9}{9}\selectfont{Winter\ Skate}}} & \multicolumn{1}{>{\raggedleft}m{\dimexpr 2in+0\tabcolsep}}{\textcolor[HTML]{000000}{\fontsize{9}{9}\selectfont{6.39}}} & \multicolumn{1}{>{\raggedleft}m{\dimexpr 2in+0\tabcolsep}}{\textcolor[HTML]{000000}{\fontsize{9}{9}\selectfont{6.05}}} \\





\multicolumn{1}{>{\raggedright}m{\dimexpr 2in+0\tabcolsep}}{\textcolor[HTML]{000000}{\fontsize{9}{9}\selectfont{Clearnose\ Skate}}} & \multicolumn{1}{>{\raggedleft}m{\dimexpr 2in+0\tabcolsep}}{\textcolor[HTML]{000000}{\fontsize{9}{9}\selectfont{5.21}}} & \multicolumn{1}{>{\raggedleft}m{\dimexpr 2in+0\tabcolsep}}{\textcolor[HTML]{000000}{\fontsize{9}{9}\selectfont{5.82}}} \\





\multicolumn{1}{>{\raggedright}m{\dimexpr 2in+0\tabcolsep}}{\textcolor[HTML]{000000}{\fontsize{9}{9}\selectfont{Red\ Hake}}} & \multicolumn{1}{>{\raggedleft}m{\dimexpr 2in+0\tabcolsep}}{\textcolor[HTML]{000000}{\fontsize{9}{9}\selectfont{8.64}}} & \multicolumn{1}{>{\raggedleft}m{\dimexpr 2in+0\tabcolsep}}{\textcolor[HTML]{000000}{\fontsize{9}{9}\selectfont{5.57}}} \\





\multicolumn{1}{>{\raggedright}m{\dimexpr 2in+0\tabcolsep}}{\textcolor[HTML]{000000}{\fontsize{9}{9}\selectfont{Smooth\ Skate*}}} & \multicolumn{1}{>{\raggedleft}m{\dimexpr 2in+0\tabcolsep}}{\textcolor[HTML]{000000}{\fontsize{9}{9}\selectfont{9.31}}} & \multicolumn{1}{>{\raggedleft}m{\dimexpr 2in+0\tabcolsep}}{\textcolor[HTML]{000000}{\fontsize{9}{9}\selectfont{5.04}}} \\





\multicolumn{1}{>{\raggedright}m{\dimexpr 2in+0\tabcolsep}}{\textcolor[HTML]{000000}{\fontsize{9}{9}\selectfont{Monkfish}}} & \multicolumn{1}{>{\raggedleft}m{\dimexpr 2in+0\tabcolsep}}{\textcolor[HTML]{000000}{\fontsize{9}{9}\selectfont{4.66}}} & \multicolumn{1}{>{\raggedleft}m{\dimexpr 2in+0\tabcolsep}}{\textcolor[HTML]{000000}{\fontsize{9}{9}\selectfont{3.23}}} \\





\multicolumn{1}{>{\raggedright}m{\dimexpr 2in+0\tabcolsep}}{\textcolor[HTML]{000000}{\fontsize{9}{9}\selectfont{Winter\ Flounder}}} & \multicolumn{1}{>{\raggedleft}m{\dimexpr 2in+0\tabcolsep}}{\textcolor[HTML]{000000}{\fontsize{9}{9}\selectfont{3.14}}} & \multicolumn{1}{>{\raggedleft}m{\dimexpr 2in+0\tabcolsep}}{\textcolor[HTML]{000000}{\fontsize{9}{9}\selectfont{3.23}}} \\





\multicolumn{1}{>{\raggedright}m{\dimexpr 2in+0\tabcolsep}}{\textcolor[HTML]{000000}{\fontsize{9}{9}\selectfont{Windowpane\ Flounder*}}} & \multicolumn{1}{>{\raggedleft}m{\dimexpr 2in+0\tabcolsep}}{\textcolor[HTML]{000000}{\fontsize{9}{9}\selectfont{3.07}}} & \multicolumn{1}{>{\raggedleft}m{\dimexpr 2in+0\tabcolsep}}{\textcolor[HTML]{000000}{\fontsize{9}{9}\selectfont{2.66}}} \\

\ascline{1.5pt}{666666}{1-3}



\end{longtable}



\arrayrulecolor[HTML]{000000}

\global\setlength{\arrayrulewidth}{\Oldarrayrulewidth}

\global\setlength{\tabcolsep}{\Oldtabcolsep}

\renewcommand*{\arraystretch}{1}

The socio-demographic conditions, and resultant vulnerabilities, of some communities may further exacerbate the impacts of offshore wind development in the Northeast such that the impacts of offshore wind development are expected to differentially \href{https://noaa-edab.github.io/catalog/wind_port.html}{impact specific coastal communities} (Fig. \ref{fig:wea-port-rev})

\begin{figure}

\includegraphics[width=1\linewidth]{../../../AppData/Local/Temp/RtmpIjNXEs/file26445de93938} \hfill{}

\caption{Percent of New England port fisheries revenue from Wind Energy Areas (WEA) from all leases (active + non-active; green), non-active leases (blue), and active leases (orange).}\label{fig:wea-port-rev}
\end{figure}

\newpage

Based on federal vessel logbook data, Little Compton, RI (13\% average and 24\% maximum) and Westport, MA (6\% average and 12\% maximum) have the highest potential revenue loss from Active Projects based on 2008-2024 total port fisheries revenue, with all other New England communities having less than 5\% (Fig. \ref{fig:wea-port-rev}). Additional fishing revenue may be lost as more areas historically used for fishing are developed for offshore wind energy. In seven Mid-Atlantic ports, New England managed species account for at least 50\% of landings from the Active Project areas by value or weight (Fig. \ref{fig:wind-rev}). Furthermore, impacts of offshore wind development may unevenly affect individual operators, with some permit holders deriving a much higher proportion of revenue from wind areas than the port-based mean.

\begin{figure}

\includegraphics[width=1\linewidth]{../../../AppData/Local/Temp/RtmpIjNXEs/file264434e3ba8} \hfill{}

\caption{Percent of New England port revenue with majority MAFMC landings from Wind Energy Areas (WEA) from  leases not unde construction (blue) and active leases (orange). }\label{fig:wind-rev}
\end{figure}

BOEM reports that cumulative offshore wind development (if all proposed projects are developed) could have moderate impacts on low-income members of vulnerable communities who work in the commercial fishing and for-hire fishing industry due to disruptions to fish populations, restrictions on navigation, and increased vessel traffic as well as existing vulnerabilities of low-income workers to economic impacts.

Top fishing communities with high \href{https://noaa-edab.github.io/catalog/engagement.html}{socio-demographic} and/or gentrification concerns such as Little Compton, RI, New Bedford, MA and New London, CT should be recognized as having potential additional vulnerability of the Active Projects and considered in decision making to reduce the social and economic impacts and aid in the resilience and adaptive capacity of these communities. In addition to fisheries landing overlaps, New Bedford, MA and New London, CT also support significant offshore wind port infrastructure needs for the Active Projects. Historically, the introduction of new industries can trigger industrial and socioeconomic gentrification of fishing ports. Competition for port space and potential pivoting of space use for offshore wind development should be monitored closely to ensure fishing communities are not adversely impacted. Additionally, offshore wind could increase recreational fishing opportunities at the turbines, potentially creating a demand for additional tourism, recreational fishing and boating port space in communities already balancing these uses with commercial fishing infrastructure, for example Point Judith, RI, and Newport, RI, and Gloucester, MA. Socio-demographic concerns also highlight communities where further resources are needed to reach underserved and underrepresented groups and create opportunities for, and directly involve, these groups in the decision-making process.

\subsubsection{Implications}\label{implications-6}

Current plans for buildout of offshore wind in a patchwork of areas spreads the impacts differentially throughout the region (Fig. \ref{fig:wind-dev-cumul}). Up to 13\% of total average revenue for major New England commercial species in lease areas could be forgone or reduced and associated effort displaced if all sites are developed. Displaced fishing effort can alter historic fishing area, timing, and method patterns, which can in turn change habitat, species (managed and protected), and fleet interactions. Several factors, including fishery regulations, fishery availability, and user conflicts affect where, when, and how fishing effort may be displaced, along with impacts to and responses of affected fish species.

Proposed wind development areas interact with the region's federal scientific surveys. Scientific surveys are impacted by offshore wind in four ways: 1. Exclusion of NOAA Fisheries' sampling platforms from the wind development area due to operational and safety limitations. 2. Impacts on the random-stratified statistical design that is the basis for scientific assessments, advice, and analyses. 3. Alteration of benthic and pelagic habitats, and airspace in and around the wind energy development, requiring new designs and methods to sample new habitats. 4. Reduced sampling productivity through navigation impacts of wind energy infrastructure on aerial and vessel survey operations.

Increased vessel transit between stations may decrease data collections that are already limited by annual days-at-sea day allocations. In the Northeast region, 14 \href{https://noaa-edab.github.io/catalog/osw_survey_impact.html}{NEFSC surveys} overlap with offshore wind development projects at varying capacities, with each of the 38 existing lease areas overlapping between 4-14 surveys. The Active Projects overlap between 10-12 surveys. Implementation of the region-wide survey mitigation program is underway with requirements to mitigate impacts to surveys included as a condition of most project approvals.

Planned development \href{https://noaa-edab.github.io/catalog/persistent_hotspots.html}{overlaps NARW} mother and calf migration corridors and a significant foraging habitat that is used throughout the year (Fig. \ref{fig:whales-wind}). Turbine presence and extraction of energy from the system could alter local oceanography and may affect right whale prey availability. For example, persistent foraging hotspots of right whales and seabirds overlap on Nantucket Shoals, where unique hydrography aggregates enhanced prey densities. Wind leases (OCS-A 0521 and OCS-A 0522) currently intersect these hotspots on the southwestern corner of Nantucket Shoals and a prominent tidal front associated with invertebrate prey swarms important to seabirds and possibly right whales. Proposed wind development areas also bring increased vessel strike risk from construction and operation vessels. In addition, there are a number of potential impacts to whales from pile driving and operational noise such as displacement, increased levels of communication masking, and elevated stress hormones.

\begin{figure}

{\centering \includegraphics[width=0.55\linewidth]{../../../AppData/Local/Temp/RtmpIjNXEs/file2644795b1796} 

}

\caption{North Atlantic Right Whale persistent hotspots (red shading) and Wind Energy Areas (black outlines).}\label{fig:whales-wind}
\end{figure}

The increase of offshore wind development can have both positive (e.g., employment opportunities) and negative (e.g., space-use conflicts) effects. Continued increase in coastal development and gentrification pressure has resulted in loss of fishing infrastructure space within ports. Understanding these existing pressures can allow for avoiding and mitigating negative impacts to our shore support industry and communities dependent on fishing. Some of the communities with the highest fisheries revenue overlap with offshore wind development areas that are also vulnerable to gentrification pressure are Point Judith and Newport, RI; and Boston and New Bedford, MA.

\paragraph{Marine Aquaculture}\label{marine-aquaculture}

Aquaculture fisheries and federally-managed fisheries could both compete or benefit each other with spatial access, shoreside infrastructure, or the supply of seafood. Unlike offshore wind, offshore aquaculture is not regulated by any federal leasing program but is permitted via the U.S. Army Corps of Engineers and the U.S. EPA. Currently, there are no federally-permitted aquaculture projects in the Northeast U.S. The marine aquaculture industry of the Northeast currently occurs in nearshore waters which are regulated by state leasing and permitting processes and federal permitting processes, as applicable. Analyses are needed to quantify the nearshore spatial distribution of aquaculture in the Northeast.

\newpage

\subsubsection{2025 Highlights}\label{highlights}

This section intends to provide a record of \href{https://noaa-edab.github.io/catalog/observation_synthesis_2025.html}{noteworthy observations reported in 2025} across the Northeast U.S. region. The full ecosystem and fisheries impacts of many of these observations are still to be determined. They should, however, be noted and considered in future analyses and management decisions.

The Northeast U.S. region experienced colder than average ocean temperatures, despite record warm \href{https://www.ncei.noaa.gov/news/global-climate-202513}{global} ocean and air temperatures. Similar to 2024, oceanographic and ecological conditions reflected cooler water and changing species abundance, distribution, and timing.

\paragraph{Northwest Atlantic Phenomena}\label{northwest-atlantic-phenomena}

The below average temperatures observed in 2024 persisted into 2025, although there are seasonal and local exceptions to this pattern. Anomalously cold surface conditions (Fig. \ref{fig:slopesea}) were recorded throughout the Northeast Shelf and were widespread across the Slope Sea for much of the year, however the waters were not as fresh as recorded in 2024. Winter bottom temperatures were also below average across much of the Northeast Shelf (Fig. 96b). Multiple oceanographic and atmospheric factors can contribute to these cooler conditions including a more southerly Gulf Stream and higher proportions of Labrador Slope and Scotian Shelf water entering the system.

\begin{figure}

{\centering \includegraphics[width=1\linewidth]{../../../AppData/Local/Temp/RtmpIjNXEs/file264467b64db7} 

}

\caption{January 2025 sea surface (left) and bottom (right) temperature differences compared to the January climatological means. Sea Surface Temperature (SST) is from the NOAA Advanced Clear-Sky Processor for Ocean (ACSPO) Super-collated SST (climatology range 2000-2020); Bottom temperature is from the GLORYS reanalysis model (climatological range 1990-2020).}\label{fig:slopesea}
\end{figure}

In 2024, the Gulf of Maine \href{https://noaa-edab.github.io/catalog/slopewater.html}{source water} entering through the Northeast Channel was near equal proportions of Warm Slope Water and cooler Labrador Slope Water (Fig. 97); data are still being processed for 2025. The colder conditions \href{https://noaa-edab.github.io/catalog/observation_synthesis_2024.html}{observed in 2024} continued into 2025 and contributed to the increased size and colder temperatures of the Mid-Atlantic \href{https://noaa-edab.github.io/catalog/cold_pool.html}{Cold Pool}.

\begin{figure}

{\centering \includegraphics[width=1\linewidth]{../../../AppData/Local/Temp/RtmpIjNXEs/file2644702c7a4b} 

}

\caption{The proportion of Warm Slope Water (WSW) and Labrador Slope Water (LSW) entering the Gulf of Maine through the Northeast Channel from 1977 to 2023. The red and teal dashed lines represent the long-term proportion averages for the WSW and LSW respectively.}\label{fig:slopewater}
\end{figure}

2025 total primary production was below average in Georges Bank and the Mid-Atlantic due to lower phytoplankton biomass and cooler sea surface temperatures. \href{https://noaa-edab.github.io/catalog/chl_pp.html?q=phyt\#chl_pp}{Phytoplankton biomass} (shown as chlorophyll a concentration) was also below average for much of 2025 (Fig. \ref{fig:chl-anom}). In particular, the winter-spring bloom period, which typically accounts for a significant proportion to total annual phytoplankton production, was shorter in duration and lower in magnitude across the entire Northeast shelf region. The fall bloom period was above average in the Gulf of Maine and Georges Bank, but near average in the Mid-Atlantic.

\begin{figure}

{\centering \includegraphics[width=1\linewidth]{../../../AppData/Local/Temp/RtmpIjNXEs/file264475726b08} 

}

\caption{Chlorophyll anomalies shown as the ratio of the monthly concentration compared to the climatological mean (1998-2020) for April (left) and October (right), 2025.}\label{fig:chl-anom}
\end{figure}

Hurricane Erin was the most notable storm in the region in 2025 and caused significant shoreside and oceanographic disturbances despite not making landfall in the Northeast. Its strong winds and large size caused mixing and weakened stratification throughout the Mid-Atlantic resulting in cooler than average surface waters across the shelf and into the Slope Sea. Along the coast, beaches from Maryland to Maine were closed due to rough surf, large waves, and rip currents, while coastal flooding led to road closures and several water rescues, particularly in New Jersey. Beach erosion was significant in some locations.

\paragraph{Northeast Shelf and Local Phenomena}\label{northeast-shelf-and-local-phenomena}

The shift to cooler waters in 2024-2025 is likely linked to multiple observations across the Northeast Shelf including the uncommon presence of Arctic zooplankton species in the Gulf of Maine, delayed migration of many species, and redistribution of some species. These shifts could affect the availability of some species to surveys or fishing, although aggregate species distributions in the cooler 2024-2025 period are tracking on the long-term trend towards northward and deeper waters (Fig. \ref{fig:species-dist}).

Mid-Atlantic scallops in the Elephant Trunk region are showing positive signs following the documented die-off. Two-year olds observed in 2024 had good survival into the 2025 survey. The Elephant Trunk region is scheduled to reopen in 2026. There was also good survival of the 2024 recruits in the southeastern Nantucket Lightship Area in 2025. In contrast, large numbers of the scallop predator \emph{Asterias vulgaris} sea stars were linked to an increased sea scallop mortality in 2024 and 2025 on Georges Bank.

Several members of the fishing community noted changes in species composition, distribution, and timing in their typical fishing grounds and attributed it to the cooler temperatures. These observations may not fully represent the entire ecosystem, but provide local context to recent events that may not be represented in other indicators. Some notable examples include:
- Several members of the fishing industry reported that it was a ``very good year'' for billfish. According to the Large Pelagic Survey, it was a record year for white marlin with more than 23,000 fish caught and released. Billfish effort may have been higher than usual due to the closure of the recreational bluefin tuna fishery in August 2025.
- Chesapeake Bay anglers reported good catch of red drum in 2024, followed by low catch and even cold stunned red drum and spotted sea trout in 2025. Scientists working with charter captains in Chesapeake Bay reported low catch rates of striped bass and red drum from June-September, but higher catch rates in the early spring and fall.
- Fishers attributed the delayed migration of black sea bass inshore, and scup migrating south for the winter using similar routes as in the early 1970s due to the cooler temperatures.
- Members of the bluefish fishery in Rhode Island reported very low landings in 2025 attributed to changes in seasonal migration path or timing.
- Some species, such as Atlantic mackerel, Illex squid and sandlance, were observed in higher abundance and wider distributions compared to recent years.
- Fishers in the Gulf of Maine and Georges Bank had mixed reports of lobsters and good catch of sea scallops.
- Fishers reported fewer warm water species in 2025 along the New Jersey coast. Others, however, noted more new species (e.g., pomapano, spadefish, triggerfish, Spanish and king mackerel) in Delaware Bay.
- Anglers also observed low spring and summer catches of gamefish in Mid-Atlantic bays and on the shelf, and high concentrations of shark species near the coast.

In Chesapeake Bay, colder than average winter 2025 temperatures (Fig. 100) were reported by state agencies as a likely cause of higher blue crab mortality rates compared to the previous winter. Colder winters generally indicate good conditions for striped bass spawning, and while the striped bass juvenile index slightly improved, it was still well below the long-term average. Several years of low striped bass recruitment is a growing concern of fisheries managers. Factors that could be influencing striped bass include below average winter-to-spring freshwater flow and above average water temperatures combined with stressful dissolved oxygen values during the summer. The continued presence of invasive blue catfish and the effect they are having on blue crab, alosines, menhaden, and striped bass populations is also a management concern.

\begin{figure}

{\centering \includegraphics[width=0.5\linewidth]{../../../AppData/Local/Temp/RtmpIjNXEs/file264452ca435d} 

}

\caption{January 2025 (left) and April 2025 (right) sea surface temperature differences compared to the respective monthly climatological means. Sea Surface Temperature (SST) is from the NOAA Advanced Clear-Sky Processor for Ocean (ACSPO) Super-collated SST (climatology range 2000-2020).}\label{fig:chesapeake-jan}
\end{figure}

\href{https://noaa-edab.github.io/catalog/ocean_acidification.html}{Ocean acidification (OA)}risk in 2025 in the Mid-Atlantic was relatively low. Compared to 2023 and 2024, there were only a few locations in 2025 where OA risk was high for Atlantic sea scallops (outer shelf, spring 2025 only), longfin squid (nearshore, summer 2025 only), and pteropods (nearshore and outer shelf, spring and summer). Gulf of Maine surface OA risk in 2025 was below 2024 levels, as indicated by aragonite saturation state (Ωa) at or near the climatological average (2006-2024).

\paragraph{Offshore Wind/Social}\label{offshore-windsocial}

Active offshore wind projects continue to be developed throughout the region. In Southern New England, South Fork Wind Farm remains the first and only commercial scale project under operation (12 turbines). Vineyard Wind 1 and Revolution Wind continued construction, and Empire Wind 1 in the New York Bight and Sunrise Wind began offshore construction in 2025 (Fig. \ref{fig:wind-dev-cumul}). Coastal Virginia Offshore Wind (CVOW) also continued construction in the Mid-Atlantic. All projects currently under construction are anticipated to be complete by the end of 2026. New London, CT and New Bedford, MA have expanded dedicated space and infrastructure for the offshore wind industry with increased port activity for the first projects under construction in southern New England. There are eight additional projects that have Construction and Operations Plan (COP) approvals (three in Southern New England and five in the Mid-Atlantic/New York Bight) that could begin construction in 2026. However, construction schedules are highly uncertain at this time.

\section{Contributors}\label{contributors}

\textbf{Editors} (NOAA NMFS Northeast Fisheries Science Center, NEFSC): Abigail Tyrell, Joseph Caracappa, Andy Beet, Brandon Beltz, Kimberly Hyde, Scott Large, Laurel Smith.

\textbf{Contributors} (NEFSC unless otherwise noted): Andrew Applegate (NEFMC), Kimberly Bastille, Heather Baertlein (NMFS Atlantic HMS Management Division), Aaron Beaver (Anchor QEA), Andy Beet, Brandon Beltz, Ruth Boettcher (Virginia Department of Game and Inland Fisheries), Mandy Bromilow (NOAA Chesapeake Bay Office), Joseph Caracappa, Samuel Chavez-Rosales, Baoshan Chen (Stony Brook University), Zhuomin Chen (UConn), Doug Christel (GARFO), Patricia Clay, Lisa Colburn, Jennifer Cudney (NMFS Atlantic HMS Management Division), Tobey Curtis (NMFS Atlantic HMS Management Division), Art Degaetano (Cornell U), Geret DePiper, Bart DiFiore (MA DMF), Emily Farr (NMFS Office of Habitat Conservation), Michael Fogarty, Paula Fratantoni, Kevin Friedland, Marjy Friedrichs (VIMS), Sarah Gaichas, Ben Galuardi (GAFRO), Avijit Gangopadhyay (School for Marine Science and Technology, University of Massachusetts Dartmouth), James Gartland (VIMS), Lori Garzio (Rutgers University), Glen Gawarkiewicz (WHOI), Maxwell Grezlik, Laura Gruenburg, Sean Hardison, Dvora Hart, Cliff Hutt (NMFS Atlantic HMS Management Division), Kimberly Hyde, Grace Jensen (WHOI), John Kocik, Steve Kress (National Audubon Society's Seabird Restoration Program), Young-Oh Kwon (Woods Hole Oceanographic Institution), Scott Large, Gabe Larouche (Cornell U), Daniel Linden, Andrew Lipsky, Sean Lucey (RWE), Don Lyons (National Audubon Society's Seabird Restoration Program), Kevin Madley, Chris Melrose, Anna Mercer, Shannon Meseck, Ryan Morse, Ray Mroch (SEFSC), Sydney Alhale (SEFSC), Brandon Muffley (MAFMC), Robert Murphy, Kimberly Murray, NEFSC staff, David Moe Nelson (NCCOS), Chris Orphanides, Stephanie Owen, Richard Pace, Debi Palka, Tom Parham (Maryland DNR), CJ Pellerin (NOAA Chesapeake Bay Office), Charles Perretti, Kristin Precoda, Grace Roskar (NMFS Office of Habitat Conservation), Jeffrey Runge (U Maine), Grace Saba (Rutgers University), Vincent Saba, Sarah Salois, Chris Schillaci (GARFO), Amy Schueller (SEFSC), Teresa Schwemmer (URI), Tarsila Seara, Dave Secor (CBL), Emily Slesinger, Angela Silva, Adrienne Silver (WHOI), Laurel Smith, Talya tenBrink (GARFO), Abigail Tyrell, Rebecca Van Hoeck, Bruce Vogt (NOAA Chesapeake Bay Office), Ron Vogel (University of Maryland Cooperative Institute for Satellite Earth System Studies and NOAA/NESDIS Center for Satellite Applications and Research), John Walden, Harvey Walsh, Sarah Weisberg, Changhua Weng, Dave Wilcox (VIMS), Timothy White (Environmental Studies Program, BOEM), Sarah Wilkin (NMFS Office of Protected Resources), Mark Wuenschel, Joseph Warren, Zhitao Yu, Qian Zhang (U Maryland).

\section{Document Orientation}\label{document-orientation}

The figure format is illustrated in Fig \ref{fig:docformat}a. Trend lines are shown when slope is significantly different from 0 at the p \textless{} 0.05 level. An orange line signifies an overall positive trend, and purple signifies a negative trend. To minimize bias introduced by small sample size, no trend is fit for \textless{} 30 year time series. Dashed lines represent mean values of time series unless the indicator is an anomaly, in which case the dashed line is equal to 0. Shaded regions indicate the past ten years. If there are no new data for 2022, the shaded region will still cover this time period. The spatial scale of indicators is either coastwide, Mid-Atlantic states (New York, New Jersey, Delaware, Maryland, Virginia, North Carolina), or at the Mid-Atlantic Bight (MAB) Ecosystem Production Unit (EPU, Fig. \ref{fig:docformat}b) level.

\begin{figure}

{\centering \includegraphics[width=1\linewidth]{C:/Users/joseph.caracappa/Documents/GitHub/READ-EDAB-SOE_reports/SOE2026_NEFMC_IR_Final_files/figure-latex/docformat-1} 

}

\caption{Document orientation. a) Key to figures: orange line - significant long-term increasing trend (purple for decreasing), dark orange line - significant increasing short-term trend (dark purple for decreasing), dashed line - time-series mean. b)The Northeast Large Marine Ecosystem: Ecological production units (EPUs) defined for the Gulf of Maine (GOM), Georges Bank (GB), and the Mid-Atlantic Bight (MAB).}\label{fig:docformat}
\end{figure}

Fish and invertebrates are aggregated into similar feeding categories (Table \ref{tab:species-groupings}) to evaluate ecosystem level trends in predators and prey.

\global\setlength{\Oldarrayrulewidth}{\arrayrulewidth}

\global\setlength{\Oldtabcolsep}{\tabcolsep}

\setlength{\tabcolsep}{2pt}

\renewcommand*{\arraystretch}{1.5}



\providecommand{\ascline}[3]{\noalign{\global\arrayrulewidth #1}\arrayrulecolor[HTML]{#2}\cline{#3}}

\begin{longtable}[c]{|p{1.00in}|p{1.00in}|p{1.00in}|p{1.00in}|p{3.00in}}

\caption{Feeding\ guilds\ and\ management\ bodies.}\label{tab:species-groupings}\\

\ascline{1.5pt}{666666}{1-5}

\multicolumn{1}{>{\raggedright}m{\dimexpr 1in+0\tabcolsep}}{\textcolor[HTML]{000000}{\fontsize{8}{8}\selectfont{Guild}}} & \multicolumn{1}{>{\raggedright}m{\dimexpr 1in+0\tabcolsep}}{\textcolor[HTML]{000000}{\fontsize{8}{8}\selectfont{MAFMC}}} & \multicolumn{1}{>{\raggedright}m{\dimexpr 1in+0\tabcolsep}}{\textcolor[HTML]{000000}{\fontsize{8}{8}\selectfont{Joint}}} & \multicolumn{1}{>{\raggedright}m{\dimexpr 1in+0\tabcolsep}}{\textcolor[HTML]{000000}{\fontsize{8}{8}\selectfont{NEFMC}}} & \multicolumn{1}{>{\raggedright}m{\dimexpr 3in+0\tabcolsep}}{\textcolor[HTML]{000000}{\fontsize{8}{8}\selectfont{State\ or\ Other}}} \\

\ascline{1.5pt}{666666}{1-5}\endfirsthead \caption[]{Feeding\ guilds\ and\ management\ bodies.}\label{tab:species-groupings}\\

\ascline{1.5pt}{666666}{1-5}

\multicolumn{1}{>{\raggedright}m{\dimexpr 1in+0\tabcolsep}}{\textcolor[HTML]{000000}{\fontsize{8}{8}\selectfont{Guild}}} & \multicolumn{1}{>{\raggedright}m{\dimexpr 1in+0\tabcolsep}}{\textcolor[HTML]{000000}{\fontsize{8}{8}\selectfont{MAFMC}}} & \multicolumn{1}{>{\raggedright}m{\dimexpr 1in+0\tabcolsep}}{\textcolor[HTML]{000000}{\fontsize{8}{8}\selectfont{Joint}}} & \multicolumn{1}{>{\raggedright}m{\dimexpr 1in+0\tabcolsep}}{\textcolor[HTML]{000000}{\fontsize{8}{8}\selectfont{NEFMC}}} & \multicolumn{1}{>{\raggedright}m{\dimexpr 3in+0\tabcolsep}}{\textcolor[HTML]{000000}{\fontsize{8}{8}\selectfont{State\ or\ Other}}} \\

\ascline{1.5pt}{666666}{1-5}\endhead



\multicolumn{1}{>{\raggedright}m{\dimexpr 1in+0\tabcolsep}}{\textcolor[HTML]{000000}{\fontsize{8}{8}\selectfont{Apex\ Predator}}} & \multicolumn{1}{>{\raggedright}m{\dimexpr 1in+0\tabcolsep}}{\textcolor[HTML]{000000}{\fontsize{8}{8}\selectfont{}}} & \multicolumn{1}{>{\raggedright}m{\dimexpr 1in+0\tabcolsep}}{\textcolor[HTML]{000000}{\fontsize{8}{8}\selectfont{}}} & \multicolumn{1}{>{\raggedright}m{\dimexpr 1in+0\tabcolsep}}{\textcolor[HTML]{000000}{\fontsize{8}{8}\selectfont{}}} & \multicolumn{1}{>{\raggedright}m{\dimexpr 3in+0\tabcolsep}}{\textcolor[HTML]{000000}{\fontsize{8}{8}\selectfont{shark\ uncl,\ swordfish,\ yellowfin\ tuna,\ bluefin\ tuna}}} \\





\multicolumn{1}{>{\raggedright}m{\dimexpr 1in+0\tabcolsep}}{\textcolor[HTML]{000000}{\fontsize{8}{8}\selectfont{Piscivore}}} & \multicolumn{1}{>{\raggedright}m{\dimexpr 1in+0\tabcolsep}}{\textcolor[HTML]{000000}{\fontsize{8}{8}\selectfont{summer\ flounder,\ bluefish,\ northern\ shortfin\ squid,\ longfin\ squid}}} & \multicolumn{1}{>{\raggedright}m{\dimexpr 1in+0\tabcolsep}}{\textcolor[HTML]{000000}{\fontsize{8}{8}\selectfont{spiny\ dogfish,\ goosefish}}} & \multicolumn{1}{>{\raggedright}m{\dimexpr 1in+0\tabcolsep}}{\textcolor[HTML]{000000}{\fontsize{8}{8}\selectfont{winter\ skate,\ clearnose\ skate,\ thorny\ skate,\ offshore\ hake,\ silver\ hake,\ atlantic\ cod,\ pollock,\ white\ hake,\ red\ hake,\ atlantic\ halibut,\ windowpane,\ acadian\ redfish}}} & \multicolumn{1}{>{\raggedright}m{\dimexpr 3in+0\tabcolsep}}{\textcolor[HTML]{000000}{\fontsize{8}{8}\selectfont{sea\ lamprey,\ sandbar\ shark,\ atlantic\ angel\ shark,\ atlantic\ torpedo,\ conger\ eel,\ spotted\ hake,\ cusk,\ fourspot\ flounder,\ john\ dory,\ atlantic\ cutlassfish,\ blue\ runner,\ striped\ bass,\ weakfish,\ sea\ raven,\ northern\ stargazer,\ banded\ rudderfish,\ atlantic\ sharpnose\ shark,\ inshore\ lizardfish,\ atlantic\ brief\ squid,\ northern\ sennet,\ king\ mackerel,\ spanish\ mackerel}}} \\





\multicolumn{1}{>{\raggedright}m{\dimexpr 1in+0\tabcolsep}}{\textcolor[HTML]{000000}{\fontsize{8}{8}\selectfont{Planktivore}}} & \multicolumn{1}{>{\raggedright}m{\dimexpr 1in+0\tabcolsep}}{\textcolor[HTML]{000000}{\fontsize{8}{8}\selectfont{atlantic\ mackerel,\ chub\ mackerel,\ butterfish}}} & \multicolumn{1}{>{\raggedright}m{\dimexpr 1in+0\tabcolsep}}{\textcolor[HTML]{000000}{\fontsize{8}{8}\selectfont{}}} & \multicolumn{1}{>{\raggedright}m{\dimexpr 1in+0\tabcolsep}}{\textcolor[HTML]{000000}{\fontsize{8}{8}\selectfont{atlantic\ herring}}} & \multicolumn{1}{>{\raggedright}m{\dimexpr 3in+0\tabcolsep}}{\textcolor[HTML]{000000}{\fontsize{8}{8}\selectfont{harvestfishes,\ smelts,\ round\ herring,\ alewife,\ blueback\ herring,\ american\ shad,\ menhaden,\ bay\ anchovy,\ striped\ anchovy,\ rainbow\ smelt,\ atlantic\ argentine,\ slender\ snipe\ eel,\ atlantic\ silverside,\ northern\ pipefish,\ atlantic\ moonfish,\ lookdown,\ blackbelly\ rosefish,\ lumpfish,\ northern\ sand\ lance,\ atlantic\ saury,\ mackerel\ scad,\ bigeye\ scad,\ round\ scad,\ rough\ scad,\ silver\ rag,\ weitzmans\ pearlsides,\ atlantic\ soft\ pout,\ sevenspine\ bay\ shrimp,\ pink\ glass\ shrimp,\ polar\ lebbeid,\ friendly\ blade\ shrimp,\ bristled\ longbeak,\ aesop\ shrimp,\ norwegian\ shrimp,\ northern\ shrimp,\ brown\ rock\ shrimp,\ atlantic\ thread\ herring,\ spanish\ sardine,\ atlantic\ bumper,\ harvestfish,\ striated\ argentine,\ silver\ anchovy}}} \\





\multicolumn{1}{>{\raggedright}m{\dimexpr 1in+0\tabcolsep}}{\textcolor[HTML]{000000}{\fontsize{8}{8}\selectfont{Benthivore}}} & \multicolumn{1}{>{\raggedright}m{\dimexpr 1in+0\tabcolsep}}{\textcolor[HTML]{000000}{\fontsize{8}{8}\selectfont{black\ sea\ bass,\ scup,\ tilefish}}} & \multicolumn{1}{>{\raggedright}m{\dimexpr 1in+0\tabcolsep}}{\textcolor[HTML]{000000}{\fontsize{8}{8}\selectfont{}}} & \multicolumn{1}{>{\raggedright}m{\dimexpr 1in+0\tabcolsep}}{\textcolor[HTML]{000000}{\fontsize{8}{8}\selectfont{barndoor\ skate,\ rosette\ skate,\ little\ skate,\ smooth\ skate,\ haddock,\ american\ plaice,\ yellowtail\ flounder,\ winter\ flounder,\ witch\ flounder,\ atlantic\ wolffish,\ ocean\ pout,\ crab,red\ deepsea}}} & \multicolumn{1}{>{\raggedright}m{\dimexpr 3in+0\tabcolsep}}{\textcolor[HTML]{000000}{\fontsize{8}{8}\selectfont{crab,unc,\ hagfish,\ porgy,red,\ sea\ bass,nk,\ atlantic\ hagfish,\ roughtail\ stingray,\ smooth\ dogfish,\ chain\ dogfish,\ bluntnose\ stingray,\ bullnose\ ray,\ southern\ stingray,\ longfin\ hake,\ fourbeard\ rockling,\ marlin-spike,\ gulf\ stream\ flounder,\ longspine\ snipefish,\ blackmouth\ bass,\ threespine\ stickleback,\ smallmouth\ flounder,\ hogchoker,\ bigeye,\ atlantic\ croaker,\ pigfish,\ northern\ kingfish,\ silver\ perch,\ spot,\ deepbody\ boarfish,\ sculpin\ uncl,\ moustache\ sculpin,\ longhorn\ sculpin,\ alligatorfish,\ grubby,\ atlantic\ seasnail,\ northern\ searobin,\ striped\ searobin,\ armored\ searobin,\ cunner,\ tautog,\ snakeblenny,\ daubed\ shanny,\ radiated\ shanny,\ red\ goatfish,\ striped\ cusk-eel,\ wolf\ eelpout,\ wrymouth,\ fawn\ cusk-eel,\ northern\ puffer,\ striped\ burrfish,\ planehead\ filefish,\ gray\ triggerfish,\ shortnose\ greeneye,\ beardfish,\ cownose\ ray,\ american\ lobster,\ cancer\ crab\ uncl,\ jonah\ crab,\ atlantic\ rock\ crab,\ blue\ crab,\ spider\ crab\ uncl,\ horseshoe\ crab,\ coarsehand\ lady\ crab,\ lady\ crab,\ northern\ stone\ crab,\ snow\ crab,\ spiny\ butterfly\ ray,\ smooth\ butterfly\ ray,\ snakefish,\ atlantic\ midshipman,\ bank\ cusk-eel,\ red\ cornetfish,\ squid\ cuttlefish\ and\ octopod\ uncl,\ spoonarm\ octopus,\ bank\ sea\ bass,\ rock\ sea\ bass,\ sand\ perch,\ cobia,\ crevalle\ jack,\ vermilion\ snapper,\ tomtate,\ jolthead\ porgy,\ saucereye\ porgy,\ whitebone\ porgy,\ knobbed\ porgy,\ sheepshead\ porgy,\ littlehead\ porgy,\ silver\ porgy,\ pinfish,\ red\ porgy,\ porgy\ and\ pinfish\ uncl,\ banded\ drum,\ southern\ kingfish,\ atlantic\ spadefish,\ leopard\ searobin,\ dusky\ flounder,\ triggerfish\ filefish\ uncl,\ blackcheek\ tonguefish,\ orange\ filefish,\ queen\ triggerfish,\ ocean\ triggerfish}}} \\





\multicolumn{1}{>{\raggedright}m{\dimexpr 1in+0\tabcolsep}}{\textcolor[HTML]{000000}{\fontsize{8}{8}\selectfont{Benthos}}} & \multicolumn{1}{>{\raggedright}m{\dimexpr 1in+0\tabcolsep}}{\textcolor[HTML]{000000}{\fontsize{8}{8}\selectfont{atlantic\ surfclam,\ ocean\ quahog}}} & \multicolumn{1}{>{\raggedright}m{\dimexpr 1in+0\tabcolsep}}{\textcolor[HTML]{000000}{\fontsize{8}{8}\selectfont{}}} & \multicolumn{1}{>{\raggedright}m{\dimexpr 1in+0\tabcolsep}}{\textcolor[HTML]{000000}{\fontsize{8}{8}\selectfont{sea\ scallop}}} & \multicolumn{1}{>{\raggedright}m{\dimexpr 3in+0\tabcolsep}}{\textcolor[HTML]{000000}{\fontsize{8}{8}\selectfont{sea\ cucumber,\ sea\ urchins,\ snails(conchs),\ sea\ urchin\ and\ sand\ dollar\ uncl,\ channeled\ whelk,\ blue\ mussel}}} \\

\ascline{1.5pt}{666666}{1-5}



\end{longtable}



\arrayrulecolor[HTML]{000000}

\global\setlength{\arrayrulewidth}{\Oldarrayrulewidth}

\global\setlength{\tabcolsep}{\Oldtabcolsep}

\renewcommand*{\arraystretch}{1}

\end{document}
